\input{preamble}



\begin{document}



\title{Seminars in Complex Geometry}

\maketitle



\phantomsection

\label{section-phantom}



\tableofcontents



\section{Everything you always wanted to know
about polygons but were too
afraid to ask}
\label{section-every-you-awlwa-wante-know}

\noindent
Alessia Mandini, UFF.
GAAG, IMPA.
September 22 and 23, 2025.

\medskip\noindent
{\bf Abstract.} Moduli spaces of polygons
form a family of Kähler
manifolds that can be constructed as a Kähler
reduction of coadjoint orbits.
These spaces have deep connections to various
areas of mathematics, including
symplectic and algebraic geometry as well as
representation theory. In this
talk, I will define these spaces and explore
their connections. I will then
discuss wall-crossing phenomena in these
spaces and demonstrate how it can be
used to determine their cohomology rings.
Finally, I will introduce the
hyperkähler analogue of these spaces, known
as hyperpolygon spaces, and describe
some of their generalizations.

\medskip\noindent



Plan of the talk
\begin{enumerate}
\item Polygons in $\mathbb{R}^3$ and
relations with
other moduli spaces.
\item Polygons in other spaces.
\end{enumerate}

Let $\alpha=(\alpha_1,\ldots,\alpha_n)\in
\mathbb{R}^n_{>0}$.
Consider $S^2$ with its usual symplectic,
Kähler form.
Also consider the product of several spheres.
There's a Hamiltonian action of
$\text{SO}(3)$ by rotations.
The moduli space of polygons is
the symplectic reduction obtained with this
Hamilatonian action,
$$
M(\alpha)=\prod
S^2_{\alpha_i}/\!/\text{SO}(3).
$$

Why is it called the space of polygons?
We think that
$$
[v_1,\ldots,v_n] \in M(\alpha) \iff
\sum_{i=1}^nv_i=0
$$
are polygons.

When smooth, $M(\alpha)$ is a
$(n-3)$-complex-dimensional
Kähler manifold. Consider
$$
\varepsilon_I(\alpha)=\sum_{i \in I}\alpha_i
-\sum_{j \in I^c}\alpha_j
$$
for some index set $I \subseteq \{1,\ldots,
n\}$.
$M(\alpha)$ is {\it smooth} if
$\varepsilon_I(\alpha)
\neq 0$ for all $I\subseteq \{1,\ldots,n\}$.
If so, we say $\alpha$ is {\it generic}.

\begin{theorem}[Kapovich-Millson]
\label{theorem-Kapovich-Millson}
$M(\alpha)$ is a complex analytic space
with (eventually) isolated singularity
(homogeneous quadratic cones).
\end{theorem}

\begin{example}
\label{examples-polygon-moduli-spaces}
\begin{enumerate}
\item $(n=3.)$ Then $M(\alpha)$ is either
empty or a point.

\item $(n=4.)$ $M(\alpha)$ is either empty or
a sphere.
\end{enumerate}
\end{example}

\begin{remark}[Hausmann-Knutson]
\label{remark-Hausmann-Knutson}
Let $M_n$ be the space of all $n$-gons
modulo rigid motions. Then $M_n$
can be equipped with a Poisson structure
for which $M(\alpha)$ are the symplectic
leaves.

\medskip\noindent
{\bf Polygons as quiver varieties.}
Consider the star-shaped quiver,
which is a distinguished point with some
points around it,
and arrows from every point to the
distinguished one.
Let the distinguished point be
$V_0=\mathbb{C}^2$ and
the rest $\mathbb{C}$. Then a representation
of this quiver
is
$$
\text{Rep}Q=\bigoplus_{i=1}^n
\Hom(\mathbb{C},\mathbb{C}^2)=\mathbb{C}^{2n}
.
$$
Now put
$$
K=(U(2) \times U(1)^n)/\Delta
$$
where $\Delta$ is the diagonal $S^1$.
This gives a Hamiltonian action on
$\mathbb{C}^{2n}$
as follows. For
$$
(A,\lambda_1,\ldots,\lambda_n)\cdot
(q_1,\ldots,q_n)
$$
we map
$q_i \mapsto  A^{-1}q_i \lambda_i$.
That is
$$
\begin{aligned}
\mu: \mathbb{C}^{2n} &\longrightarrow K^* \\
(q_1,\ldots,q_n) &\longmapsto
\left(\sum_{n=1}^n (q_iq_i^*
),\ldots,\frac{1}{2}|q_i|^2\ldots\right)
\end{aligned}
$$

\begin{theorem}[Hausmann-Knutson]
\label{theorem-Hausmann-Knutson}
$$
\mathbb{C}^{2n}/\!/_{(0,\alpha)}K=M(\alpha)
$$
\end{theorem}

\begin{proof}
Uses
$$
\begin{aligned}
\mathbb{R}^3  &\xrightarrow{\simeq}
\mathfrak{su}(2)^* \\
v_i &\longmapsto (q_iq_i^*)_0
\end{aligned}
$$
\end{proof}
\end{remark}

$$
\xymatrix{
&\mathbb{C}^{2n}\ar[dl]_{/\!/_\alpha U(1)^n}
\ar[dr]^{/\!/U(2)}\\
\prod S^2 \ar[dr]_{/\!/_\alpha \text{SO}(3)}
 &  &
 \text{Gr}(2,\mathbb{C}^n)\ar[dl]^{/\!/
 _\alpha
 U(1)^n}\\
&M(\alpha).
}
$$
So we have (I think former symplectic
reduction)
$$
\mu_{U(1)^n}:\text{Gr}(2,\mathbb{C}^n)\to
\mathbb{R}^n
$$
\medskip\noindent
{\bf Walls and wall crossing.}
The walls are
$$
W_I=\{\alpha \in
\mathbb{R}^n_{>0}:\varepsilon_I(\alpha)=0\}
$$
for $I \subseteq \{1,\ldots,n\}$.

To understand wall crossing suppose we have
a wall $W_I$, with $a^c$ in the wall,
$\alpha^+$ on one side and $\alpha^-$ on the
other.

Not that
$$
\varepsilon_I(\alpha^+)>0 \iff
\sum_{k \in I}\alpha_i^+>
\sum_{ j \in I^c}\alpha_j^+
$$
Define
$$
\begin{aligned}
M_{I^c}(\alpha^+)
&=\{(v_1,\ldots,v_n) \in M(\alpha^+):
v_i=\lambda v_j \forall i,j \in I^c,
\lambda>0\}\\
&=M(\tilde{\alpha}),\qquad \tilde{\alpha}=
\left(\alpha_{i_1},\ldots,\alpha_{i_k},
\sum_{j
\in I^c}\alpha_j\right)
\end{aligned}
$$
Then $\varepsilon_I(\alpha^+)<0$
implies $M_I(\alpha^-) \subseteq
M(\alpha^-)$.

$$
\xymatrix{
&\tilde{M} \subseteq E=M_I \times
M_{I^c}\ar[dd]^{\text{blow up}}\\
M(\alpha^+)\ar[dr]&  &  M(\alpha^-)\ar[dl]\\
&M(\alpha^c)
}
$$
\begin{remark}
\label{remark-modul-space-polyg-propo}
\begin{itemize}
\item $M(\alpha)\cong M(\sigma(\alpha))$
for $\sigma$ a permutation on the order of
sets.
\item $M(\alpha)$ is conformally
symplectomorphic
to $M(\lambda \alpha)$ for all $\lambda>0$.
\end{itemize}
\end{remark}

\begin{definition}
\label{definition-short-and-long}
Let $\alpha \in \mathbb{R}^n_{>0}$.
$I \subseteq \{1,\ldots,n\}$ is {\it short}
if $\varepsilon_I(\alpha)<0$ and {\it long}
if  $\varepsilon_I(\alpha)>0$.
\end{definition}

\begin{example}
\label{example-wall-crossing}
We did an example of wall crossing.
The relevant manifolds $M_{I^c}(\alpha)$ and
$M_I(\alpha)$
were projective spaces.
\end{example}

\medskip\noindent
Now recall our first quotient
$$
\xymatrix{
\mu_{U(1)}^{-1}(\alpha)\subseteq
\text{Gr}(2,\mathbb{C}^n)
\ar[d]^{/\!/U(1)^n}\\
M(\alpha)
}
$$
Let $c_i$ be the first Chern
classes associated to the
$n$ $S^1$-bundles above
(by taking reduction in stages).

\begin{theorem}[Haussmann-Knutson,M.]
\label{theorem-Haussmann-Knutson-M}
The $c_i$ generate the cohomology of
$M(\alpha)$.
\end{theorem}

\begin{theorem}[Guillemin-Stenberg]
\label{theorem-Guillemin-Stenberg}
In this setup, for $\alpha$ generic,
$$
H(M)=\mathbb{C}[c_1,\ldots,c_n]/\text{Ann}
(\text{Vol}(M(\alpha)))
$$
i.e., $Q(c_1,\ldots,c_n) \in
\text{Ann}(\text{Vol}M(\alpha))$
$\iff$ $Q\left(\frac{\partial }{\partial
\alpha_i},\ldots,
\frac{\partial }{\partial \alpha_1}\right)
\text{Vol}(M(\alpha))=0$
\end{theorem}

\begin{theorem}[Takakura, The Koi]
\label{theorem-Takakura-The-Koi}
$$
\text{Vol}(M(\alpha))
=-\frac{(2\pi)^{n-3}}{(n-3)!}
\sum_{I\text{
long}}(-1)^{n-|I|}\varepsilon_I(\alpha)^{n-3}
.
$$
\end{theorem}

\begin{example}
\label{example-Delta1}
According to Example
\ref{example-wall-crossing} (which I
did not copy) we find that
$\alpha_1 \in \Delta_1$ gives
$\text{Vol}(M(\alpha_1))=2\pi^3(1-2\alpha_3)
^2$
and
$$
H(M(\alpha_1))=\mathbb{C}[c_3]/(c_3^3)
$$
\end{example}

\medskip\noindent
{\bf Polygon game.}
Let $G$ be a Lie group and $\mathfrak{g}$ its
Lie algebra.
Then the co-adjoint orbits
$\mathfrak{g}^*  \supseteq
\mathcal{O}_{\xi_i}$
are symplectic manifolds with the KKS form.
Then
$$
\begin{aligned}
G\text{ acts on }\Pi \mathcal{O}_{\xi_i}
&\longrightarrow \mathfrak{g}^* \\
(A_i,\ldots\alpha_n) &\longmapsto \sum A_i
\end{aligned}
$$
Then the quotient
$$
M(\xi)=\Pi \mathcal{O}_{\xi_i}/\!/_0 G
$$
generalizes the previous construction,
which we get with $G=\text{SU}(2)$.

It could be interesting to investigate
which of the next constructions can be
generalized
to other Lie groups.

\begin{theorem}[Sotillo-Floentins-Gadihl]
\label{theorem-Sotillo}
Wall-crossing for $\text{SU}(m)$.
\end{theorem}

\medskip\noindent
{\bf Bending action.}
We consider a polygon of $n$ sides
and put some diagonals that don't intersect.
We introduce some
notion of ``bending'' that allows
to define a Hamiltonian function.
In turn, this defines a torus action on
$M(\alpha)$
and a moment map.

For $n=2$ we obtain the moment map
$$
\begin{aligned}
\mu: M(\alpha) &\longrightarrow \mathbb{R}^2
\\
p &\longmapsto (\ell_1(p),\ell(p))
\end{aligned}
$$
and we can see the moment polytope in
$\mathbb{R}^2$.

\begin{remark}
\label{remark-nonva-nonin-diago}
Any system of $(n-3)$ non-vanishing
and non-intersecting diagonals determine
a torus action of $T^{n-3}$ on $M(\alpha)$.
\end{remark}

\medskip\noindent
{\bf Relations of polygon spaces to other
moduli spaces.}
Representations of the fundamental group
of the punctured sphere in $SU(2)$, i.e.
$$
\text{Rep}(\pi_1(S^2\setminus
\{p_1,\ldots,p_n\},\text{SU}(2))/\text{SU}(2)
\cong
M(\alpha)
$$
This can be put very explicitly:
$$
\begin{aligned}
&\text{Rep}(\pi_1(S^2\setminus
\{p_1,\ldots,p_n\},\text{SU}(2))\\
&=\{(g_1,\ldots,g_n)\in
\text{SU}(2)^n:g_1\cdot\ldots\cdot
g_n=\text{Id},
t_rg_i=2\cos\pi \alpha_i\}
\end{aligned}
$$
See [Agapito-Godinho]. That's all we
will say about this example.

\medskip\noindent
Now consider the moduli space of parabolic
bundles.
Consider a holomorphic bundle of rank 2 over
$\mathbb{C}P^{2}$.
Choose some points $D=\{x_1,\ldots,x_n\}$ in
$\mathbb{C}P^{2}$
and a flag
$E_{x_i}=E^{x_i,1}\supseteq
E^{x_i,2}\supseteq \{0\}$
for every $x_i \in D$. Each of the
$E_{x_i,j}$ is isomorphic to $\mathbb{C}$
(1-dimensional).
A {\it quasi-parabolic bundle}
is an holomorphic bundle $E$
with such a choice of flag.
A {\it parabolic bundle} is a quasi-parabolic
bundle with a choice of parabolic weights
$0< \beta_1(x_i)< \beta_2(x_i)<1$.

There is also a notion of stability:
$$
\begin{aligned}
\text{pdeg}E&=\text{deg}E+ \sum_{i=1}^n
(\beta_1(x_i)+\beta_2(x_i))\\
\mu(E)&=\frac{\text{pdeg}(E)}{\text{rank}(E)}
\end{aligned}
$$
We say $E$ is {\it (semi)stable} if
$µ(E)> \mu(L)$ ($\mu(E) \geq \mu(L)$)
for any $L \subseteq E$ parabolic
subbundle.

Then we obtain that
$\mathcal{M}_{\pi,d}(\beta)$
is the moduli space of (semi)-stable
parabolic bundles on $\mathbb{P}^1$
of rank $r$ and degree $d$.
In particular $\mathcal{M}_{2,0}(\beta)$
is the moduli space of (semi)-stable
parabolic
bundles on $\mathbb{P}^1$ of rank 2 and
degree 0
holomorphically trivial.

\begin{theorem}[Jeffrey, Godinho, M.]
\label{theorem-Jeffrey-Godinho}
For generic $ \alpha$,
$M(\alpha)$ is diffeomorphic to
$\mathcal{M}_{2,0}(\beta)$
whenever
$\alpha_i=\beta_2(x_i)-\beta_1(x_i)$.
\end{theorem}

\begin{proof}[Idea of proof]
The correspondence can be made quite explicit
as a map
$$
\begin{aligned}
M(\alpha) &\longrightarrow
\mathcal{M}_{2,0}(\beta) \\
[q_1,\ldots,q_n] &\longmapsto
\substack{E=\mathbb{C}P^{1}\times\mathbb{C}^2
 \\ E_{x_i}\supset E_{x_i,1}\supset
 X_{x_i,2}\supset \{0\}}\end{aligned}
$$

\end{proof}

\medskip\noindent
Now consider the quiver variety we described
above:
$$
\text{Rep}Q=\bigoplus_{i=1}^n
\Hom(\mathbb{C},\mathbb{C}^2)=\mathbb{C}^{2n}
.
$$
Now put
$$
K=(U(2) \times U(1)^n)/\Delta
$$
where $\Delta$ is the diagonal $S^1$.

But this time put
$$
\text{Rep}\tilde{Q}=\bigoplus_{i=1}^n
\Hom(\mathbb{C},\mathbb{C}^2)
\oplus \bigoplus_{i=1}^n
\Hom(\mathbb{C}^2,\mathbb{C})
$$
We also have an action of $K$, and we get
$$
K\text{ acts on }\text{Rep}\tilde{Q} = T^*
\mathbb{C}^{2n},
$$
a so-called {\it hyperHamiltonian action}.
The quotient
$$
X(\alpha)=/\!/\!/_{\substack{(0,\alpha) \\
(0,0)}}K=
\frac{\mu_{\mathbb{R}}1(0,\alpha) \wedge
\mu^{-1}_{\mathbb{C}}(0,0)}{K}
$$
called {\it hyperpolygon space}.

Here
$$
\begin{aligned}
\mu_\mathbb{R}:T^*\mathbb{C}^{2n}&\to
\mathcal{K}^* \\
\mu_{\mathbb{C}}:T^*\mathbb{C}^{2n}&  \to
\mathcal{K}_{\mathbb{C}}^* \\
(q_1,\ldots,q_n,p_1,\ldots,p_n)&  \mapsto
(\sum_{i=1}^n(p_i,q_i)_0,\ldots)
\end{aligned}
$$
It turns out that $X(\alpha)$ is smooth
if and only if $\varepsilon_1(\alpha)=0$ for
all
$I \subseteq \{ 1,\ldots,n\}$.
When smooth,  $X(\alpha)$ is a hyperkähler
manifold
(non-compact, unfortunately),
$M(\alpha)=\{[0_{1i},0] \in X(\alpha)\}$.

\begin{theorem}[Boalch]
\label{theorem-Boalch}
$X(\alpha)$ is the moduli space of polygons
for $\text{GL}(2,\mathbb{C})$.
\end{theorem}

\medskip\noindent
{\bf Parabolic Higgs bundles.}
Let $E$ be a parabolic bundle as before, i.e.
$E \in \mathcal{M}_{0,2}(\beta)$.
A {\it Higgs field on $E$} is
$$
\phi \in H^{0}(\mathbb{P}^1,\text{SPEnd}(E)
\otimes K_{\mathbb{P}^1}(D)))
$$
where a {\it strongly parabolic endomorphism}
is  $f:E \to E$
such that $f(E_{x_i})\subseteq E_{x_{i+1}}$
for all $i$.
(Fix details!)

A {\it parabolic Higgs bundle} is (probably
the pair $(E,\phi)$).
A parabolic Higgs bundle $(E,\phi)$ is {\it
(semi)stable}
if  $\mu(E) > \mu(L)$ ($\mu(E) \geq \mu(L)$)
for all
$L$ parabolic Higgs subbundle.

Then $\mathcal{N}_{r,a}^{0,1}(\beta)$ is the
moduli
space of parabolic Higgs bundles over
$\mathbb{P}^1$
with rank $r$ and degree $d$ (with fixed
determinant,
and traceless; two notions that we will not
define here).


\begin{theorem}[Goldinho, M.; Biswar,
Florentino, Godinho, M.]
\label{theorem-GBSFG}
Let $\alpha$ be generic, let $\beta$
be such that
$\alpha_i=\beta_2(x_i)-\beta_1(x_i)$.
Then the hyperpolygon space $X(\alpha)$
is symplectomorphic to the moduli space
of parabolic HIggs bundles over
$\mathbb{P}^1$, $\beta$-stable,
holomorphically trivial (with fixed
determinant and trace-free).


\begin{proof}[Idea of proof]
We can define a correspondence
$$
\begin{aligned}
X(\alpha) &\longrightarrow \mathcal{H}(\beta)
\\
[p,q] &\longmapsto
\substack{E=\mathbb{C}P^{1}\times
\mathbb{C}^2
 \\ E_{x_i}\simeq E_{x_i,1}\supset
 E_{x_i,2\supset \{0\}}}
\end{aligned}
$$
No we can use the moment map condition
that the sum of the residues is zero,
where $\text{Res}_{x_i}\phi=(p_i,q_i)$
where $\phi$ is the unique function
that satisfies the last equality by
the residue theorem.
\end{proof}
\end{theorem}

\medskip\noindent
{\bf Wall-crossing for moduli spaces of
parabolic bundles.}
See ``Translation''from Thaddeus.
What happens is that there is an $S^1$-action
on $X(\alpha)$ given by
$$
\lambda\cdot[p,q]=[\lambda p,q].
$$
(This can be paralleled with the $S^1$-action
on $\mathcal{H}(\beta)$,
given by
$\lambda\cdot[E,\phi]=[E,\lambda\phi]$.)

This action has a moment map
$$
\mu_{S^1}([p,q])=\frac{1}{2}\sum_{i}|p_i|^2
$$
Fixed points are [Konno]
$$
X_S=\{[p,q] \in X(\alpha):
S,S^c \text{ are straight, }p_j=0 \forall j
\in S\}
$$
for all $|S|\geq 2$ short. Here $S\subseteq
\{1,\ldots,n\}$
is {\it straight} if $q_i=\lambda_j q_j$
for $\lambda_j>0$ $\forall  i,j \in S$
(which in the case of polygons means the
edges are aligned).

Let $U_S$ be the flow down from $X_S$
Then, crossing a wall  $W_I$ as defined
above,
i.e. $W_I=\{\alpha \in
\mathbb{R}^n_{>0}:\varepsilon_S(\alpha)=0\}$,
``replaces'' $U_S$ by $U_{S^c}$.

Let $\alpha$ and $\tilde{\alpha}$ be generic,
then $X(\alpha)$ is diffeomorphic to
$X(\tilde{\alpha})$.
The wall crossing for $M(\alpha)$
involves
$$
\begin{aligned}
M_S(\alpha^+)&=U_S \cap M(\alpha^+)\\
M_{S^c}(\alpha^-)&=U_{S^c}\cap M(\alpha)
\end{aligned}
$$
This is thought of as a Mukai transform.

\medskip\noindent
{\bf Generalization.}
See [Florentino, Godinho, Sotillo] for wall
crossing.
See also [Fisher] PhD thesis, [Fisher,
Rayan].
[Hausel et al], [Rayan-Schopnick].

\section{Vertex algebras and special holonomy
on quadratic Lie algebras}
\label{section-verte-algeb-speci-holon}

\noindent
Mario García Fernandez, ICMAT.
GAAG, IMPA.
September 24, 2025.

\medskip\noindent
{\bf Abstract.}
The chiral de Rham complex (CDR) is a sheaf
of vertex algebras on any smooth
manifold, introduced by Malikov, Schechtman
and Vaintrob, which provides a
formal quantization of the non-linear sigma
model in mathematical physics.
Motivated by the algebra of chiral symmetries
in two-dimensional superconformal
field theories, vertex algebra embeddings on
the CDR have been studied for
special holonomy Riemannian manifolds, thanks
mainly to the work of Heluani,
Zabzine, and collaborators, with interesting
applications to the elliptic genus.
In these lectures, we will discuss extensions
of some of these results to the
case of special holonomy manifolds with
skew-torsion. The presence of torsion
typically allows for continuous symmetries in
the geometry, with an enhanced
interplay with Lie theory and algebra, as
well as the application of techniques
from generalized geometry.


\medskip\noindent
Based on joint work with Luis Álvarez Cónsul,
Andoni De Arriba de la Hera. arXiv:2012.01851
(IMRN '24).

\medskip\noindent
{\bf Motivation.}
Let $(M^n,g)$ be a Riemannian spin manifold
with
parallel spinor $\nabla^g \varphi=0$.
Then $\text{hol}(g) \subset G_\varphi
\subseteq \text{SO}(n)$,
and $\text{Ric}(g)=0$.

There is a construction by Markov-S-V that
puts a sheaf
of vertex algebras $\mathcal{V} \to M^n$.
How to construct special embeddings of
trivial vertex algebras
in the cohomology of $\mathcal{V}$, i.e.
$\mathcal{V}_\varphi \hookrightarrow
H^*(\mathcal{V})$
by Heluani et al.

\medskip\noindent
Applications. Construction of topological
invariants.
Elliptic genus, [Borisov-L].
How to understand mirror symmetry using
vertex algebras [Borisov].
Holography [Witten].

\medskip\noindent
{\bf Geometry in algebra.}
We shall do geometry in quadratic Lie
algebras.
Recall that a {\it quadratic Lie algebra} is
 $(\mathfrak{g},[\cdot,\cdot],
 (\cdot,\cdot))$ where
$(\mathfrak{g},[\cdot,\cdot])$ is a Lie
algebra
(in this course we use $\mathbb{R}$ as base
field)
and $(\cdot,\cdot)$ is a symmetric bilinear
invariant form.

\begin{example}
\label{example-quadratic-Lie-algebra}
Let $R$ be a Lie algebra with
$\left\langle{}\cdot,\cdot\right\rangle{}_R: R \otimes R \to
\mathbb{R}$.
Take $\mathfrak{g}=R \oplus R^*$,
and pick $H \in \Lambda^{3}R^*$.
Put $H(a,b,c)=\left\langle{}[a,b],c\right\rangle{}$, and
define
$$
[v+\alpha,w+\varphi]=[v,w]-\varphi([v,-])+
\alpha([v,-])+H(v,w,-)
$$
for $v,w \in R$ and $\varphi,\alpha \in R
^*$. This turns out to be a bracket.

If you like geometry you can pick $K$
compact, $\text{Lie}K=R$,
this is ``Generalized geometry on $T K \oplus
T^* K$''
à la Hitchin.
\end{example}

\medskip\noindent
QLA:
\begin{itemize}
\item Courant algebroids /$\{*\}$.
\item  Symplectic supermanifolds.
\end{itemize}

\medskip\noindent
\begin{definition}
\label{definition-generalizaed}
\begin{enumerate}
\item A {\it (generalized) metric} on
$(\mathfrak{g},(\cdot,\cdot))$ is
$G \in \text{End}(\mathfrak{g})$ with
$G^2=\text{Id}$, $(G,G)=(\cdot,\cdot)$,

This gives $\mathfrak{g}=V_+ \oplus V_-$
where $G$ acts as identity on the
first term and as $-\text{Id}$ on the second
one.
$(G_-,-)$ is a non-degenerate pairing.
$(-,-)_{V_I}$ non-degenerate tensor.

\item A divergence $\varphi \in \mathfrak{g}
^*$.

\item A {\it connection} is
$$
D: \mathfrak{g} \to \mathfrak{g}^* \otimes
\mathfrak{g}
$$
such that $(D_ab,c)+(b,D_ac)=0$.
So that $D \in \mathfrak{g}^*  \otimes
\Lambda^{2}\mathfrak{g}$
(where we identify $\mathfrak{g}$ with its
dual using the pairing).

$D$ is  {\it compatible} with $G$ if
$[D_a,G]=0$. This condition says that $D$
splits into four operators:
$$
\xymatrix{
D^+_+&  &  D_+^-\\
&  D\ar[ul]\ar[ur]\ar[dl]\ar[dr]\\
D^+_-& & D_-^-
}
$$
\item  Let $a \in \mathfrak{g}$, $a=a_++a_-$.
A {\it generalized connection} satisfies
$D_a^+=[a_-,b_+]_+ +[a_+,b_-]_-$. (This will
happen to things
living in $\mathcal{D}^0$, see below.)
\end{enumerate}
\end{definition}

\begin{definition}
\label{definition-}
Given $D$ define
\begin{enumerate}
\item $\varphi_D(a)=-\text{tr}Da$.
\item $\text{Torsion}(AX):T_D \in
\Lambda^{2}\mathfrak{g}^*$,
$$
T_D(a,b,c)=(D_ab-D_ba-[a,b],c)+(D_ca,b).
$$
(Just copying the definition of torsion.)
\end{enumerate}
\end{definition}

\begin{lemma}
\label{lemma-}
Define
$\mathcal{D}^0(G,\varphi)=\{D:G\text{-
compatible
}, \varphi_D=\varphi\}$.
Then this is non-empty, but is not a point.
Furthermore, $\forall  D \in
\mathcal{D}^0(G,\varphi)$,
$$
D_{a_-}b_+=[a_-,b_+]_+,\qquad
D_{a_+}b_-=[a_+,b_-]_-.
$$
\end{lemma}

(Checking non-emptyness is just writing out
the equations.)

\medskip\noindent
Given $G$ consider $Cl (V_+)$,
 $a_+\cdot a_+=(a_+,a_+)$ for $a_+ \in V_+$.

Fix an irreducible representation for
$Cl(V_+)$: $S_+$.
Given  $D \in \mathcal{D}^0(G,\varphi)$:
\begin{enumerate}
\item $D_-^{\delta_+} \in V_-^* \otimes
\text{End}((S_+)\supset
V_-^*\otimes \Lambda^{2}V_+ \ni D_-^+$.
\item An operator like Dirac operator:
$$
\begin{aligned}
\underline{D}^+: S_+ &\longrightarrow S_+ \\
\zeta &\longmapsto \sum_{j=1}^{\dim
V}e_j^+\cdot D_{e_j^+}\zeta.
\end{aligned}
$$
\end{enumerate}

\begin{lemma}
\label{lemma-independendant}
$D^{S_+}$ and $\underline{D}^+$ are
independnt of
$D \in \mathcal{D}^0(G,\varphi)$.
\end{lemma}

\begin{definition}
\label{definition-Killing-spinor-equations}
$(G,\varphi,\zeta)$, $\varphi \in S_+$
satisfies
{\it Killing spinor equations} if
$$
KSE\qquad
\underbrace{D_-^{S_+}\zeta=0,}
_{\text{Gravitino
Eq.}}\qquad
\underbrace{\underline{D}^+\varphi=0
}_{\text{Dilatino Eq.}}$$
\end{definition}

Expectation.
\begin{enumerate}
\item
$\mathcal{M}=\{(G,\varphi,\zeta):KSE\}/
\mathfrak{g}$
special metric.
\item Given a solution $(G,\varphi,\zeta)$
then
$$
V_{(g,\varphi,\zeta)}\hookrightarrow
V^k(\mathfrak{g}).
$$
(Where I think $V^k(\mathfrak{g})$ is the
Kac-Moody affinization.)
\end{enumerate}

\begin{remark}
\label{remark-Lie-algebra}
In the case $T K \oplus T^*K: \mathcal{M}$
2-stack. See [Bursztyn].
\end{remark}

\begin{lemma}
\label{lemma-Ricci-tensors}
Associated to $(G,\varphi)$ there are
well-defined ``Ricci tensors''
$\text{Ric}^+ \in V_- \otimes V_+$,
$\text{Ric}^- \in V_+ \otimes V_-$.
$$
\text{Ric}^+(a_-,b_+)=\text{Tr}(C_+ \to
R_D(c_+,a_-)b_+)
$$
$$
R_D(a,b)=[D_a,D_b]_c-D_{[a,b]}c,\qquad D \in
\mathcal{D}^0(G,\varphi).
$$
\begin{remark}
\label{remark-in-geometric-setup}
In geometric setup, the vanishing of the
Ricci corresponds to
the motion equations of some physical
supersymmetry theory.
\end{remark}

\begin{proposition}
\label{proposition-KSE}
If $(G,\varphi,\zeta)$ is solution of KSE,
then
 $\text{Ric}^+_{G,\phi}=0$.
If $[\varphi,G]=0 \implies
\text{Ric}^-_{G,\varphi}=0$.
\end{proposition}
\end{lemma}

\begin{proof}
$\text{Ric}^+(a-,-)\cdot
\varphi=[\underline{D}^+,D_{a_-}^{S_+}\zeta
-D_{D_+a_-}^{S_+}\zeta=0$.
$[\varphi,G]=0 \implies
\text{Ric}^+(a_-,b_+)=\text{Ric}^-(b_+,a_-)$.
\end{proof}

\medskip\noindent
{\bf Generalized Ricci flow.}
Finding solutions to these equations.
Evaluating by
$$
G_t^{-1}\partial_tG_t=-2(\text{Ric}^+-
\text{Ric}^-)
$$
$\Hom(V_+,V_-) \oplus \Hom(V_-,V_+)$.
$G G+G - G=0$,
$G=\text{Id}_{V_+}-\text{Id}_{V_-}$.

\begin{exercise}
\label{exercise-}
\begin{enumerate}
\item Prove STE for GRF.
\item Assuming there exists a solution of
KSE, prove long
time existence and convergence. (See theorem
by Streets, Jordan, GF;
and a book by Streets, GF.)
\end{enumerate}
\end{exercise}

\begin{remark}
\label{remark-sigma-model}
For physicists the generalized Ricci flow
(GRF) is the GRF of a $2d$
$\sigma$-model with target a compact Lie
group
$K$,  $\mathfrak{g}=R \oplus R^*$
\end{remark}

Let's explain something in this situation.
consider a compact Lie group $K$ and
$\mathfrak{g}=R \oplus R^*$.
A generalized metric: $\mathfrak{g}=V_+
\oplus V_-$,
$(-,-)|_{V_\pm}$ non-degenerate,
$(-,-)|_{V_+}>0$.
$V_+=\mathcal{C}^b\{X+g(X)\}$,
$X \in R$, $g \in S^2(R^*)$, $b \in
\Lambda^{2}R^*$,
$H=H_0+\overline{\partial}$.

\medskip\noindent
{\bf Case $\dim V_+=2n$.}
Complex pure spinor $\varphi$ on $V_+$ is
equivalent $V_+ \otimes \mathbb{C}=\ell
\oplus \overline{\ell}$,
where $\ell=\{0 \in V_+ \oplus \mathbb{C}:0
\cdot\zeta=0\}$.
$(\ell,\ell)=0$,
$(\overline{\ell},\overline{\ell})=0$.

This gives a complex structure
$J_\zeta:V_+ \to V_+$, $J_\zeta=i
\text{Id}_\ell-i
\text{Id}_{\overline{\ell}}$.

\begin{lemma}
\label{lemma-solutions-KSE}
$(G,\varphi,\zeta)$ with $\zeta$ pure
satisfies KSE if and only if
\begin{enumerate}
\item $[\ell,\ell]\subset \ell$.
\item Take bars
$\{\varepsilon_j,\overline{\varepsilon}_j\}
^n_{j=0}$,
$\varepsilon_j \in \ell$,
$\overline{\varepsilon}_j\in \overline{\ell}$
$(\varepsilon_j,\overline{\varepsilon}_j) =
s_{j,k}$.
$\sum_{j=1}^n[\varepsilon_j,
\overline{\varepsilon}_j=-J_\zeta
\varphi_+$.
\end{enumerate}
\end{lemma}
(This sum is a moment map condition.)

Consider $\mathcal{L}=\{\ell \subset
\mathfrak{g} \otimes \mathbb{C}:
\dim\ell=n,\ell\cap\overline{0},(-,-)|_{\ell
\oplus \overline{\ell}}
\text{non-degenerate}\}$.
This is a complex manifold.
$T_\ell \mathcal{L}
\cong \Hom(\ell,\overline{\ell}\oplus V_-
\oplus \mathbb{C})$,
$V_- \otimes \mathbb{C}=(\ell \oplus
\overline{\ell})^\perp$
Pick a vector $\dot \ell = \dot J + \dot G
\in
\Hom(\ell,\overline{\ell}\oplus V_- \oplus
\mathbb{C})$
That is, $\dot J \in \Hom(\ell,
\overline{\ell})$,
$\dot J :V_+ \to V_+$,
$\dot J J+ J \dot J=0$. (Deformations of the
complex structure.)
$G \in \Hom(V_+,V_-) \oplus \Hom(V_-,V_+)$.

\begin{proposition}[Romero]
\label{proposition-Romero}
$\mathcal{L}$ has a pseudo-Kähler structure
preserved by the $\mathfrak{g}$-action and
there exists
a moment map
$$
\begin{aligned}
\mu: \mathcal{L} &\longrightarrow
\mathfrak{g}^* \\
\ell &\longmapsto
\frac{i}{2}\sum_{j=1}^n[\varepsilon_j,
\overline{\varepsilon}_j],-),
\end{aligned}
$$
which is the quantity we mentioned in Lemma
above.
\end{proposition}

\begin{proof}
Use the natural complex structure on the
space of complex structures
(found by Fujiki),
$$
\text{tr}_{V_+}(J \dot G_2,\dot
G_1)-\text{tr}_{V_+}(J \dot J_{\varepsilon}
\dot J_1).
$$
\end{proof}

Problem: prove that
$\mathcal{M}=\{(G,\varphi,\zeta):\zeta\text{
puse}\}/\mathfrak{g}$
is a pseudo-Kähler manifold.

\begin{enumerate}
\item Pseudo-Kähler.
\item Shifter symplectic stuff.
\end{enumerate}

\begin{example}
\label{example-SU2xSU1}
Take $K=\text{SU}(2)\times\text{U}(1)\cong
S^3\times S^1$
and $\mathfrak{g}=R \oplus R^*$.
Take generators $v_1,v_2,v_3$ of
$\text{SU}(2)$ and $v_4$ of $\text{U}(1)$.
We have $[v_2,v_3]=-v_1$, $[v_3,v_1]=-v_2$,
$[v_1,v_2]=-v_3$.
Put $H_\ell=\ell v^{123}$. $\ell \in
\mathbb{R}$. $x,a \in \mathbb{R}_{>0}$.
$$
g_{x,a}=\frac{a}{x}\left(\sum_{i=1}
^3\omega^{\otimes
2}
+x^2 (v^4)^{\otimes 2}\right)
$$
$$
V_+=\{ x+g_{xa}\}\subset \mathfrak{g}
$$
$$
I_X v_4=xv_1,\qquad U_X v_2=v_3, \qquad
\varphi=-xv_4
$$
\end{example}

\begin{exercise}
\label{exercise-solution-of-KSE}
If $\ell=\frac{a}{x} \implies $ solution of
KSE.
\end{exercise}

Furthermore $[\varphi_+,\ell] \subset \ell$
(holomorphic divergence).

The $a$ parameter is naturally complexified
by  $b=yv^{23}$.
$y+ia=z$.
$$
V=\text{log}
\left(\frac{\omega^2_{x,n}}{v^{1234}}\right)=
\text{log}a.
$$
KSV hyperbolic metric, metric on
$\mathbb{H}$.

\medskip\noindent
{\bf $\dim V_+=7$}.
A real spinor on $V_+$ is equivalent to $\phi
\in (\Lambda^{3}V_+^* )_{>0}$.
The action of $\text{GL}(\mathbb{R}^7)$ on
$\Lambda^{3}(V_+^*)$.
The space of spinors $S_+=\mathbb{R}^8
=\mathbb{R}^7\oplus\mathbb{R}
\left\langle{}\varphi\right\rangle{}$.
$\phi(x,y,z)=\left\langle{}x\cdot y\cdot z\cdot
\zeta,\zeta\right\rangle{}$.

\begin{remark}
\label{remark-7}
$v^1,\ldots,v^7 \cong \mathbb{R}^6 \oplus
\mathbb{R}$.
$\phi=(v^{12}+v^{34}+v^{56})\wedge v^7+
\text{Re}((v^1+i
v^2)\wedge(v^3+iv^4)\wedge(v^5+iv^6))$.
\end{remark}

\begin{example}
\label{example-SU2+SU2+R}
Take $R=\text{SU}(2)\oplus \text{SU}(2)\oplus
\mathbb{R}$.
$\mathfrak{g}=R \oplus R^*$.
$e,s \in \mathbb{R}$, $H=sv^{123}+\ell
v^{456}$.
$\phi=\omega \wedge \zeta+\Omega^+$,
$\zeta=\sqrt{\varepsilon/\ell}v^7$,
$\omega=\sqrt{s\ell}(v^{14}+v^{25}-v^{36})$,
$\Omega^+=\sqrt{s^3}v^{123}+\ell
\sqrt{s}v^{156}-\ell \sqrt{s}v^{345}$.
\end{example}

\section{Non-Kähler Hodge Lefschetz theory
and the Bianchi identity}
\label{section-non-kahle-hodge-lefsc-theor}

\noindent
Arpan Saha, UNICAMP.
Geometric Structures Seminar, IMPA.
September 25, 2025.

\medskip\noindent
{\bf Abstract.}

Being Kähler imposes severe constraints on
the cohomology of compact complex
manifolds such as the Hard Lefschetz
property, and the question of how far this
generalises beyond the class of Kähler
manifolds has been of great interest for
a while. In this talk, I shall report on
ongoing joint work with Mario García
Fernández and Raúl González Molina that
abstracts out the definition of a
variation of Hodge--Lefschetz structure and
provides evidence that, under
certain natural assumptions, such a structure
exists more generally on
distinguished subspaces within moduli spaces
of Bismut--Ricci-flat metrics that
are pluriclosed up to source terms. In
particular, these distinguished subspaces
may be regarded as replacements for the
Kähler cone, with affine structure
modelled on a subspace of the (1,1) Aeppli
cohomology of the compact complex
manifold.

\medskip\noindent
{\bf Broad motivation.}
As we move on the parameter space associated
to a CY manifold we encounter walls.
The Kähler cone is contained in this
parameter space, and we may cross
its walls. When we cross a wall
we obtain a birational transformation
(I think this means that the moduli
on either side are birational).
But the Kähler condition is not a
birational invariant.

\medskip\noindent
{\bf Hodge-Lefschetz theory.}
For $(X,J,\omega)$ Kähler,
the cohomology $H^{\bullet}(X)$ comes with:
\begin{itemize}
\item the Lefschetz operator $L$,
given by $[\omega]\wedge-$.

\item Hodge star operator.
\item $\Lambda=* L *$ satisfying
the $\mathfrak{sl}_2$ relations:
$$
[L,\Lambda]=H,\qquad [H,L]=2L,\qquad
[H,\Lambda]=-2\Lambda.
$$
These relations are equivalent to
the Hard Lefschetz property that
$$
L^q:H^{d-q} \to H^{d+q}.
$$
\item The Poincaré pairing
(which is a different pairing from Lefschetz)
of the form
$$
H^{d-q}\cong (H^{d+q})
$$
given by integration.

\item Hodge-Riemann bilinear relations
which is a positive definite form given by
$\int_X\cdot \wedge *\cdot$.

\end{itemize}

\noindent
Moving about on the Kähler cone gives
a variation of Hodge-Lefschetz structure.

\begin{definition}
\label{definition-varia-hodge-lefsc-struc}
A {\it variation of Hodge Lefschetz structure
(VHLS)}
of weight $d$ over a manifold $K$ consists of
\begin{itemize}
\item a graded real vector bundle
$E=\bigoplus_{q=0}^d E^q \to K$
with $\text{rk}_\mathbb{R}E^0=1$.

\item An isomorphism $E^0 \otimes TK
\xrightarrow{\simeq}E^1$
(this is reminiscent of the usual definition
of variations of Hodge structures, and
we will not use it in this talk).

\item Endomorphism fields
$L, \Lambda,H \in
\mathcal{A}^0(K,\text{End}(E))$
satisfying the $\mathfrak{sl}_2$ relations
that
$$
[L,\Lambda]=H,\qquad [H,L]=2L,\qquad
[H,\Lambda]=-2\Lambda.
$$
and that
$$
H|_{E_q}=(2q-d)|_{\text{id}_{E^q}}.
$$
\item An involution $* \in
\mathcal{A}^0(K,\text{End}(E))$
such that $* L *=\Lambda$.

\item A nondegenerate symmetric pairing
$P \in \mathcal{A}^0(\mathcal{K},E^* \otimes
E^*)$
w.r.t. $L$ and $*$ are self-adjoint.

\item (like the Gauss-Manin connection)
a flat connection $D:
\mathcal{A}^0(K,E)\to\mathcal{A}^1(K,E)$
such that $DH=0=DP$.
\end{itemize}
The VHLS is said to be {\it positive} if
$P(*-,-)$
is positive definite.
\end{definition}

\noindent
{\bf Calabi's dream beyond Kähler manifolds.}
Let $(X,\Omega)$ be a compact Calabi-Yau
manifold,
where $\Omega$ is the volume form.
In general such a manifold is not Kähler.

\begin{example}[Non-Kähler manifolds]
\label{example-non-kahler}
These manifolds are not Kähler:
Heisenberg algebra quotiented by some
lattice,
torus bundle, Hopf surface.
\end{example}

\noindent
By [Yau], the Kähler cone is the moduli
space of Kähler Ricci-flat metrics.

We look for PDEs on  $X$ such that
\begin{itemize}
\item Assumption C (Calabi). Moduli
space of alutions is an open set $K$
of an affine space modelled on
a subspace $\mathbb{V}$ of
$(1,1)$-Aeppli cohomology group,
which is given by
$$
H^{1,1}_A(X)=\frac{\Ker
\partial\overline{\partial}|_{\Omega^{1,1}}}
{\text{Im}
\partial+\text{Im}\overline{\partial}}
$$

\item Assumtion D (dilaton).
A volume function $v:K\to \mathbb{R}_{>0}$
such that $\underline{d}\text{log} V$
is nowhere zero (where $\underline{d}$
denotes
Aeppli differential), and the bilinear form
$g_K:=-Dd \text{log}V$
is nondegenerate (basically the Hessian
metric).
(Ideally this would be positive definite.)

\item Assumption V (vector field).
$Z=g_K^{-1}d \text{log}V$
is such that $DZ$ is invertible
(when thought of as an endomorphism)
$g_K(Z,Z)=d \in \mathbb{Z}_{>0}$,
 $$
(D_{(DZ)^{-1}})^{d+1}V=0,
$$
\end{itemize}

\begin{remark}
\label{remark-pairing}
$$
(D_{(DZ)^{-1}})^dV=\int_X \wedge \wedge
\wedge.
$$
\end{remark}

\begin{theorem}[García Fernández-González
Molina-AS]
\label{theorem-GGA}
If assumptions C, D and V all hold with  $d
\leq 3$,
then $K$ admits a VHLS of weight $d$.
\end{theorem}

\begin{proof}[Idea of proof for $d=3$]
$$
E=\underbrace{\mathbb{R}_K}_{E^0} \oplus
\underbrace{TK}_{E^1}
\oplus \underbrace{T^*K}_{E^2} \oplus
\underbrace{\mathbb{R}_K^*}_{E^3}.
$$
Define $L$ in each of the terms of the direct
sum as
$$
L1=Z,\qquad Lv=D^2_vV,\qquad
L_\alpha=i_Z\alpha
$$
and $*$ as
$$
*1=V 1^\vee,\qquad  *v=VD^2_0
\text{log}V=Vg_K(v,-).
$$
\end{proof}

\medskip\noindent
Recall the Hull-Strominger
system, which is given by two conditions
on the connection and the metric
$$
d\left(\|\Omega\|_\omega
\frac{\omega^{n-1}}{(n-1)!}\right)=0,\qquad
F_\theta \wedge
\frac{\omega^{n-1}}{(n-1)!}=0.
$$
This determines $d
d^c\omega=\alpha\left\langle{}F_\theta \wedge
F_\theta\right\rangle{}$.

\section{Formality of Dolbeault cohomology}
\label{section-forma-dolbe-cohom}

\noindent
Misha Verbitsky, IMPA.
Geometric Structures Seminar, IMPA.
October 29, 2025.

\medskip\noindent
{\bf Abstract.}
A quasi-isomorphism of differential graded
algebras (DGA) is a multiplicative
map inducing an isomorphism on cohomology. A
DGA is called  formal if it can be
connected by a chain of  quasi-isomorphisms
to its cohomology algebra. For the
de Rham algebra of a compact manifold,
formality is an important  topological
property which has many implications (such as
vanishing of Massey products). For
the Dolbeault DGA of a compact complex
manifold, the situation is not well
understood, unless it is Kahler. The
Dolbeault DGA of a torus, symmetric spaces
and hyperkahler manifolds are known to be
formal.  We prove that the Dolbeault
DGA of a complex nilmanifold is formal only
if it is a torus, and the Dolbeault
algebra of (0,p)-forms is formal if and only
if the complex structure is
abelian. This is a joint work with Tommaso
Sferruzza.

\medskip\noindent
Upshot: the Dolbeault dga of a nilmanifold
$M$ is never formal.

\medskip\noindent
If a space is formal, Massey products vanish.
Unfortunately, Massey products are not
well-defined.
Let $a,b,c \in \Lambda^*(M)$ be closed forms
such that their cohomology classes satisfy
$[a][b]=[b][c]=0$.
Let  $\alpha,\gamma \in \Lambda^*(M)$
be such that $d(\alpha)=a \wedge b$ and
$d(\gamma)b \wedge c$.
Then the {\it Massey product} of $a,b,c$
is the class  $[\alpha \wedge
c-(-1)^{\tilde{a}}a \wedge \gamma]$
which is only defined up to
 $\left\langle{}\text{img}L_a + \text{img}
 L_c\right\rangle{}$
where $L_a$ and $L_c$ are multiplication by
$a$ and $c$ operators.

The {\it Heisenberg nilmanifold} is the
quotient of
the Heisenberg group $G$ of upper triangular
matrices
with $1$ on the diagonal over
$G_{\mathbb{Z}}$
of the same group with integer entries.
The Heisenberg nilmanifold fibers over a
torus
with fiber a circle:
$$
\xymatrix{
0\ar[r]&S^1\ar[r]&G/G_{\mathbb{Z}}\ar[r]&
\mathbb{R}^2/\mathbb{Z}^2\ar[r]&0
}
$$
which in fact comes from the central
extension
$$
\xymatrix{
0\ar[r]&\mathbb{R}\ar[r]&G\ar[r]&\mathbb{R}
^2\ar[r]&0.
}
$$


\begin{lemma}
\label{lemma-heise-nilma-masse-produ}
Massey products on  $G/G_{\mathbb{Z}}$ are
non-zero.
\end{lemma}

\begin{proof}
Step 1. There is an action of $G$ in the
cohomology of $G$.
It's not hard to see that the cohomology
classes on $G/G_{\mathbb{Z}}$
can be represented by right $G$-invariant
forms,
and the cohomology of $G/G_{\mathbb{Z}}$ is
equal to the cohomology of the complex
of right-$G$-invariant forms on $G$.

\medskip\noindent
Step 2. In fact, this complex is the same as
the
Chevalley-Eilenberg complex of the Lie
algebra $\mathfrak{g}$ of $G$,
and the cohomology is the Lie algebra
cohomology.

\medskip\noindent
Step 3. From a basis $a,b,t$ of
$\mathfrak{g}$ with the only nontrivial
commutator $[a,b]=t$ and $\alpha,\beta,\tau$
dual basis in $\mathfrak{g}^*$
with the only nontrivial differential
$d\tau=\alpha\wedge\tau$,
we obtain a basis $\alpha \wedge \beta,
\alpha \wedge \tau, \beta \wedge \tau$
in $\Lambda^2(\mathfrak{g}^*)$ with
$d|_{\Lambda^2\mathfrak{g}^*}=0$.
This gives $\text{rk}
H^1(G/G_{\mathbb{Z}})=2$
and $\text{rk}H^2(G/G_{\mathbb{Z}})=2$.

\medskip\noindent
Step 4. Compute the Massey product of
$\alpha,\beta,\tau$:
it does not vanish.
\end{proof}

\medskip\noindent
Digression: this is Sullivan's motivation for
working on dgas.
Consider the class of spaces whose homotopy
and homology
groups are vectors spaces over $\mathbb{Q}$
(and it's better to also assume that $\pi_1$
is nilpotent).
There is a equivalence of categories between
old category
and homotopy category of rational homotopy
spaces.
And in fact it's also equivalent to the
(homotopy?)
category of dga algebras.

\medskip\noindent
Now we move to differential graded algebras
(dga).
A morphism of dgas (which I think only means
that commutes with
the differentials) is called {\it
quasi-isomorphism} if it
induces isomorphisms on all $H^i(A^\bullet)$.

A dga is {\it formal} if it is connected to
$H^\bullet(A^\bullet)$
by a chain of quasi-isomorphisms.

Recall the $dd^c$-lemma. It states that for
the twisted
differential $d^c=Jd J^{-1}$, if a form
$\alpha$ is $d$-closed and $d^c$-exact,
(or $\partial$-exact and
$\overline{\partial}$-exact),
then $\alpha\in \text{img}dd^c$.

\begin{lemma}
\label{lemma-quasi-isomorphisms}
The map $(\Lambda^\bullet(M)_{d^c},d)\to
(\Lambda^\bullet(M),d)$
is a quasi-isomorphism.
\end{lemma}

\begin{proof}
Uses Hodge decomposition of a form to prove
induced maps
on cohomology are injective. Then prove
surjectivity.
\end{proof}

\begin{lemma}
\label{lemma-kahler-formal}
A compact Kähler manifold is formal (i.e. the
usual cohomology algebra
is formal).
\end{lemma}

\begin{theorem}
\label{theorem-kahler-dolbeault-formal}
$M$ Kähler then
$(\Lambda^\bullet(M),\overline{\partial})$ is
formal.
\end{theorem}

\medskip\noindent
A manifold $M$ is {\it nilmanifold} if it has
an action of aa nilpotent Lie
group. All nilmanifolds are obtained as
qutient spaces $M=G/\Gamma$
for a cocompact lattice in a nilpotent Lie
group.

I gather there are two things:
one is an {\it integrable complex structure
on a Lie algebra}, which
is a subalgebra
$\mathfrak{g}^{1,0}\subset \mathfrak{g}
\otimes_\mathbb{R} \mathbb{C}$
such that $\mathfrak{g}^{1,0} \oplus
\overline{g^{1,0}}
=\mathfrak{g} \otimes \mathbb{C}$.
Integrability is
$[\mathfrak{g}^{1,0},\mathfrak{g}^{1,0}]
\subset \mathfrak{g}^{1,0}$.

And then there's a {\it complex nilmanifold},
which is $M=G/\Gamma$
with a complex structure $I$ such that $G$
has
a right-invariant complex structure and $G
\to M$ is holomorphic.
In this case we say $I$ is {\it invariant}.
Not all complex structures on a nilmanifold
are invariant:
a compact torus admits a non-Kähler complex
structure that is not invariant.
A Kodaira surface, a locally trivial
holomorphic fibration over $T$ with fiber
$T'$ with non-trivial Chern class, for $T$
and $T'$ elliptic curves,
is a nilmanifold.

Recall the Chevalley-Eilenberg cohomology
on $\Lambda^\bullet\mathfrak{g}^*$ where the
Lie bracket becomes
a differential by Cartan formula.
More generally,

\begin{proposition}
\label{proposition-cheva-eilen-diffe}
Let $w \in \Hom(\Lambda^2V,V)$.
Consider the dual map $d_w:V^* \to
\Lambda^2V^*$.
Extend this map to $d_w:\Lambda^kV^*  \to
\Lambda^{k+1}V^*$
using Leinbiz rule. Then $d^2_w=0$ if and
only if $w$
defines the Lie algebra structure on $V$.
\end{proposition}

\noindent
We may identify the Chevalley-Eilenberg
complex of the Lie algebra  $\mathfrak{g}$
with the
complex of left-invariant differential forms
in its
Lie group $G$. This defines a natural
homomorphism of dgas
$(\Lambda^\bullet\mathfrak{g}^*,d) \to
(\Lambda^\bullet(G/\Gamma),d)$
for any discrete group $\Gamma \subset G$.

\begin{theorem}[Nomizu]
\label{theorem-nomizu}
Let $M=G/\Gamma$ be a nilmanifold.
Then the map above is a quasi-isomorphism.
\end{theorem}

As a corollary, we see that nilmanifolds,
with exception of tori, are never formal.

\begin{proposition}
\label{proposition-}
Let $(\mathfrak{g},I$) be a complex structure
on a Lie algebra. Then we may consider the
Dolbeault
cohomology of
$\partial_w:\Lambda^\bullet(\mathfrak{g}^*)
\to\Lambda^2(\mathfrak{g}^*)$.
Then $[\cdot,\cdot]_{\partial_w}$ is defined
as follows: let
$x,y \in \mathfrak{g}^{1,0}$, $x',y' \in
\mathfrak{g}^{0,1}$.
Then
$$
\begin{aligned}
[x,y]_{\partial_w}&=[x,y]\\
[x,x']_{\partial_w}&=[x,x']^{1,0}\\
[x',y']_{\partial_w}&=0.
\end{aligned}
$$
\end{proposition}

\noindent
This implies that $(CE,\overline{\partial})$
is formal if and only if
$[\cdot,\cdot]=0$.

\medskip\noindent
We intend to apply the construction made
before for the usual
Lie algebra cohomology but for Dolbeault
cohomology.
Unfortunately there is no analogue of
Nomizu's theorem Theorem
\ref{theorem-nomizu}.

\noindent
[Missing slides]

\medskip\noindent
We say a dga has {\it Poincaré duality} if
the product
map $H^i(A^* ) \times H^{n-i}(A^*)\to
H^n(A^*)=\mathbb{C}$
is a non-degenerate pairing in cohomology.
A morphism of dgas $(B^*,d)\to(A^*,d)$, where
$(B^*,d)$ admits
Poincaré duality, that induces injective maps
on cohomology
is called a {\it domination}.

\begin{theorem}[Milivojevic,Stelzig,Zoller
arXiv
\href{http://arxiv.org/abs/2306.12364}%
{2306.12364}]
\label{theorem-msz}
Let $\psi:(B^*,d)\to (A^*,d)$ be a domination
of dga.
If $(B^*,d)$ is not formal, then $(A^*,d)$ is
not formal.
\end{theorem}

\begin{theorem}[Barberis-Dotti-V., 2007]
\label{theorem-bdv7}
A complex nilmanifold has trivial canonical
bundle,
trivialized by an invariant holomorphic top
form.
\end{theorem}

As corollaries we get that the natural dga
morphism
$(\Lambda^\bullet\mathfrak{g}^*,
\overline{\partial}_w)\to
(\Lambda^\bullet
G/\Gamma,\overline{\partial})$ is a
domination,
and that the Dolbeault dga of a nilmanifold
$M$ is never formal.

\section{Dolbeault cohomology of
nilmanifolds}
\label{section-dolbe-cohom-nilma}

\noindent
Misha Verbitsky, IMPA.
Geometric Structures Seminar, IMPA.
November 6, 2025.

\medskip\noindent
{\bf Abstract.}
A nilmanifold is a (left) quotient of a
nilpotent Lie group by a cocompact
lattice. If this Lie group is equipped with a
left-invariant complex structure,
the quotient is called a complex nilmanifold.
Examples of complex nilmanifolds
include Kodaira-Thurston surface and many
other non-Kahler complex manifolds.
The differential graded algebra of invariant
differential forms on a nilmanifold
has the same cohomology as a nilmanifold, as
shown by Nomizu. I am going to
explain a Dolbeault cohomology version of
this theorem. A Lie algebra of a
nilmanifold admits a natural rational
structure, induced by the cocompact
lattice. If the complex structure operator,
acting on the Lie algebra, has
rational coefficients, Console and Fino
proved that the Dolbeault cohomology of
a nilmanifold is equal to the Dolbeault
cohomology of its invariant forms. I
will describe an attempt to prove this result
for a general nilpotent Lie
algebra, based on number-theoretic arguments.

\medskip\noindent
Let $G$ be a nilpotent Lie group.
[For me this would mean that the Lie algebra
of $G$ is
nilpotent as in Lie Algebras Definition
\ref{lie-algebras-definition-nilpo-%
solva-lie-algeb},
but according to Misha this would be the case
for $G$
connected only; so there's must some other
notion
of nilpotent Lie groups.]
A {\it nilmanifold} $M$ is a manifold with a
transitive action
of a nilpotent Lie group. This means that
$M=G/\Gamma$.

\begin{theorem}[]
\label{theorem-}

\end{theorem}

\noindent
Let $G$ be a real Lie group
and $I \in \text{End}(TG)$ with
$I^2=-\text{Id}$.
Suppose $I$ is right-invariant.
Then $I$ is integrable iff
$[T^{1,0}G,T^{1,0}G]\subset T^{1,0}G$
iff
$\mathfrak{g}^{1,0}\subset\mathfrak{g}
\otimes\mathbb{C}$
where $\mathfrak{g}^{1,0}$ is the
$\sqrt{-1}$-eigenspace of $I$.

A {\it complex nilmanifold} is
$(G,I)/\Gamma$.

\begin{example}
\label{example-iwasawa-manifold}
An {\it Iwasawa manifold} is a quotient of
the set
of upper triangular matrices with 1 on the
diagonal
(this group is $U_3(\mathbb{C})$) by a
lattice $\Gamma$.
For example, we can take
$\Gamma=U_3(\mathbb{Z}[\sqrt{-1}])$.
\end{example}

\noindent
The complex structure $I$ is bi-invariant iff
it is paralelizable.

\medskip\noindent
Let $E$ be an elliptic curve with an ample
line bundle $L$.
We put an action on $\text{Tot}L$ given by
multiplication by
$\lambda \in \mathbb{C}$ for $|\lambda|>0$.

\begin{theorem}
\label{theorem-de-rham-cheva-eilen}
De Rham cohomology of a nilmanifold coincides
with the
CE cohomology of $\mathfrak{g}$.
\end{theorem}

\begin{theorem}[Console-Fino]
\label{theorem-console-fino}
For complex nilmanifolds,
$H^{\bullet,\bullet}(\mathfrak{g}^\bullet,
\overline{\partial}$
and
$H^{\bullet,\bullet}_{\overline{\partial}}(M)
$
are isomorphic.
\end{theorem}

\noindent
There's several results addressing a converse
situation
(when the algebras are isomorphic) and other
situations
involving fibrations.

\medskip\noindent
Now we make a sudden change to deformation
theory.

\begin{lemma}
\label{lemma-}
Complex spaces on $V=\mathbb{R}^{2n}$ are
uniquely determined
by $W \in V \otimes_{\mathbb{R}}\mathbb{C}$
such that $W \oplus \overline{ W}=V \otimes
\mathbb{C}$
which is an element in $\text{Gr}(n, V
\otimes \mathbb{C})$.
\end{lemma}

\begin{definition}
\label{definition-}
Let $\mathfrak{g}$ be a Lie algebra and
$\mathfrak{g}_\mathbb{C}$
its complexification.
The set of complex structures on
$\mathfrak{g}$ is $\text{Comp}(\mathfrak{g})$
which is $\{ W \in
\text{Gr}_0(\mathfrak{g}_\mathbb{C}:
W \subset W \subset
\mathfrak{g}_{\mathbb{C}}\text{ is a
subalgebra}\}$
since being a subalgebra means
$[\mathfrak{g}^{1,0},\mathfrak{g}^{1,0}]
\subset \mathfrak{g}^{1,0}$.
\end{definition}

We can define a variety out of this:
$$
\overline{\text{Comp}(\mathfrak{g}
_{\mathbb{C}}}
=\{W \in
\text{Gr}(n,\mathfrak{g}_{\mathbb{C}}:
\text{$W$ subalgebra}\}
$$
is a complex subvareity of
$\text{Gr}(n,\mathfrak{g}_{\mathbb{C}})$.


\begin{definition}
\label{definition-field-of-defition}
Let $X \subset \mathbb{C}P^n$ be a projective
subvariety.
Consider the smallest field $k \supset
\mathbb{Q}$ such that
any ideal $I$ of $X$ [ideal in the coordinate
algebra of $X$?]
can be generated by polynomials with
coefficients
in $k[k^n]$ is called a {\it field of
definition}.
\end{definition}

\noindent
Apparently there's some discussion about this
notion
not being easily defined.

\begin{lemma}
\label{lemma-}
$\text{Gal}_{\mathbb{Q}}(k)$ acts on
$\mathbb{P}_k^n$.
Let $X \subset \mathbb{P}^n$ be a projective
subvariety
and $\Gamma \subset
\text{Gal}_{\mathbb{Q}}(k)$
the stabilizer of $X$.
Then the field of definition of $X=\{\lambda
\in k:\Gamma(\lambda)=\lambda\}$.
\end{lemma}

\noindent
We have a simpler definition of field of
definition for a point
(which is apparently the former definition
when the variety is a point)

\begin{definition}
\label{definition-field-of-definition-point}
For a point $(z_1,\ldots,z_n) \in
\mathbb{C}^n$, its
{\it field of definition} is the field
generated by $z_1,\ldots,z_n$.
\end{definition}

\noindent
So, for instance, a rational point has
transcendence degree 0.

\begin{remark}
\label{remark-}
Let $X$ be an $n$-dimensional algebraic
variety
defined over $\mathbb{Q}$ or any finite
extension of $\mathbb{Q}$,
and let $\pi:X\to \mathbb{C}^n$ be finite
dominant over $\mathbb{Q}$.
\end{remark}


\noindent
We have a fun theorem:

\begin{theorem}
\label{theorem-fun-theorem}
Let $X$ be an irreducible complex algebraic
variety defined over $\mathbb{Q}$.
Then the group
$\text{Aut}_\mathbb{Q}(\mathbb{C})$ acts
transitively
on points of maximal transcendence degree in
$X$.
\end{theorem}

\begin{definition}
\label{definition-algebraic-point}
We say that a point $x \in X$ is {\it
algebraic} if
its transcendence degree over $k_X$ is zero.
\end{definition}

\noindent
It's not very clear whether the variety $X$
is projective or affine.

\begin{lemma}
\label{lemma-algebraic-points-are-dense}
The set of algebraic points in
$X\subset\mathbb{C}P^n$ is dense
in the analytic topology.
\end{lemma}

\begin{theorem}
\label{theorem-aut-orbit-of-point}
Let $X$ be a complex (affine) algebraic
variety defined over $\mathbb{Q}$.
Consider an
$\text{Aut}_{\mathbb{Q}}\mathbb{C}$-orbit of
a point
$p \in X$. Then the closure of its orbit is a
minimal subvariety $Z \subset X$ defined
over $\mathbb{Q}$ and containing $p$.
\end{theorem}

\noindent
Now we introduce a Galois action on Dolbeault
cohomology.

\begin{theorem}
\label{theorem-}
$\mathfrak{g}$ nilpotent Lie algebra, $M$
corresponding nilmanifold,
$\sigma \in
\text{Aut}_{\mathbb{Q}}(\mathbb{C})$,
and let $J=\sigma(I) \in \text{Comp}(G)$.
Then $\dim H^{p,q}_{\overline{\partial}}(M,I)
=\dim H^{p,q}_{\overline{\partial}}(M,J)$ for
all $p,q$.
\end{theorem}

From this it follows that

\begin{lemma}
\label{lemma-console-fino-complex-structure}
For a nilmanifold with complex structure $I$,
if the field of definition of $I$ is
$\mathbb{Q}$,
then the Console-Fino map
$H_{\overline{\partial}}^{p,q}(\mathfrak{g})
\to
H_{\overline{\partial}}^{p,q}(M)$ is an
isomorphism.
\end{lemma}

\noindent
The idea is that if we know the result for
algebraic
points then we'll have it for transcendental
points as well.
Namely,

\begin{lemma}
\label{lemma-if-conso-fino-algeb-then-hold}
If Console-Fino is an isomorphism for all
algebraic complex structures,
then it holds for all transcendental complex
structures too.
\end{lemma}

\begin{proof}
Let $I \in \text{Comp}(\mathfrak{g})$ be
trascendental.
Consider the variety
$\overline{\text{Aut}_{\mathbb{Q}}\mathbb{C}}
$,
which is algebraic and thus contains
algebraic points.
Then Console-Fino is an isomorphism on such
points.
Then it is so in a neighbourhood of this
point.
Then it is so in the orbit of this
neighbourhood.
\end{proof}

\noindent
Consider $(\mathfrak{g},I)$ with $I$ with
number field $k$.
Consider
$\tilde{\mathfrak{g}}=\bigoplus_{\sigma \in
\text{Gal}(k:\mathbb{Q})}
(\mathfrak{g},\sigma(I))$…
this constructions leads to a Lie algebra
$(\tilde{\mathfrak{g}},\tilde{I})$
equipped with a complex structure defined
over $\mathbb{Q}$.
This satisfies Console-Fino. We can lift
$(\tilde{M},\tilde{I})$
since it remains Dolbeault harmonic --- which
is the main tool used by
Colsole-Fino in their theorem.
This is how to proof Console-Fino for complex
structures over a closed field.

\section{Holomorphic symmetric differentials,
semiample bundles,
fundamental groups}
\label{section-holom-symme-diffe-semia}

\noindent
Ernesto Mistretta, Università di Padova.
Algebra Seminar, IMPA.
November 12, 2025.

\medskip\noindent
{\bf Abstract.}
We will present some recent results and some
open conjectures relating
holomorphic symmetric differentials and the
geometry or topology of a compact
complex manifold, we will illustrate the
difference in the projective, or
Kaehler, or compact complex cases. Then we
will apply some elementary
constructions on semiample vector bundles to
partially answer some of these
questions.



\medskip\noindent



\noindent
As a motivation,
$$
\xymatrix{
&\text{differential forms}\ar[ld]\ar[rd]\\
\text{topology}&&\text{geometry}
}
$$
$X$ compact complex manifold, $E$
(holomorphic) vector bundle.
$L$ line bundle, the  {\it Kodaira-Itaka
dimension} is
$\text{KI}(X,L)=k$ if $h^0(X,L^{\otimes
m})=\mathcal{O}(m^k)$.

{\it Holomorphic $p$-differentials} are
elements of
$H^0(X,\text{Sym}^m\Omega^p_X)$ for some
$m>0$.

A conjecture by Mumford: suppose that $X$ is
a compact
Kähler manifold, then $X$ is rationally
connected if and only if
for every  $m>0$ and for every $p>0$,
$H^0(X,\text{Sym}^m\Omega^p_X)=0$.

Some relevant results are the following.

\begin{theorem}[Campana-Demailly-Peterwell,
2015]
\label{theorem-cdp}
Suppose $X$ is a smooth and projective
variety over $\mathbb{C}$.
$X$ is rationally connected if and only if
for every
ample line bundle there exists a positive
constant $C_A>0$,
such that $\forall p>0$, $\forall K>0$ and
$\forall m>C_A K$,
$$
H^0(X,\text{Sym}^m\Omega^p_X \otimes A^K)=0.
$$
\end{theorem}

\noindent
Recall that if $X$ is a Kähler manifold and
rationally connected,
then $X$ is projective and $\pi_1(X)=0$.

\begin{theorem}[Brunarbe-Campana, 2016]
\label{theorem-bc16}
Let $X$ be compact Kähler.
Suppose that $\forall m>0$, $\forall p>0$,
$H^0(X,\text{Sym}^m\Omega^p_X)=0$.
Then $X$ is projective and $\pi_1(X)=0$.
\end{theorem}

\noindent
What if we restrict to $p=1$?
$H^0(X,\text{Sym}^m\Omega^1_X)$?

A conjecture by Esnault says: let $X$ be a
compact Kähler manifold.
Suppose that $\forall m>0$,
$H^0(X,\text{Sym}^m\Omega^1_X)=0$.
Then $|\pi_1(X)|<\infty$.

\begin{theorem}[Burnabbe-Klinger-Totsaro]
\label{theorem-bkt}
Suppose $X$ is a compact Kähler manifold.
Suppose that there exists a
finite-dimensional
representation of $\rho:\pi_1(X)\to
\text{GL}_r(k)$
over some field $k$ such that
$|\pi(\pi_1(X))|=\infty$.
Then there exists $m>0$ such that
 $H^0(X,\text{Sym}^m\Omega_X^1) \neq 0$.
\end{theorem}

\noindent
Notice the relationship between Theorem
\ref{theorem-bkt}
and the conjecture by Esnault.

\medskip\noindent
Conjecturally, we have:
\begin{enumerate}
\item No holomorphic symmetric
$p$-differentials for all $p$
if and only if $X$ is rationally connected.
\item No holomorphic symmetric
1-differentials, them $|\pi_1(X)|<\infty$.
\end{enumerate}

\noindent
Question. Suppose there are some (many)
nonzero holomorphic symmetric
1-differentials, does this imply that
$|\pi_1|=\infty$?
What if $X$ is compact Kähler? And if $X$ is
projective?

Let's first answer the easiest of these
questions.
Let $X$ be compact Kähler with
$H^0(X,\Omega^1_X)\neq 0$.
By Hodge theorem,
$b_1(X)=\dim_\mathbb{Q}H^1(X,\mathbb{Q})
=2\dim_\mathbb{C}H^0(X,\Omega_X^1)$.
Also,
$H_1(X,\mathbb{Z})=\text{Ab}(\pi_1(X))$.
In particular, if $b_1 \neq 0$,
then
$H_1(X,\mathbb{Z})=\mathbb{Z}^{b_1}\oplus$
torsion.

\begin{theorem}[Brotbeck-Drondeau,2018]
\label{theorem-brotbeck-drondeau18}
$X \subseteq \mathbb{P}^N$ very gen. complete
intersection
of high degree such that $\dim X\leq N/2$.
Then $\Omega_X^1$ is ample.
(In particular, $\text{Sym}^m\Omega_X^1$
is globally generated for $\gg m$.)
Now if $\dim X\geq 2$, then
$\pi_1(X)=\pi_1(\mathbb{P}^N)=0$.
\end{theorem}

\noindent
In this case $\omega_X=\det(\Omega_X^1)$ is
an ample line bundle.
So the Kodaira dimension of $X$ is
$\text{kod}(X)=\text{kod}(X,\omega_X)=\dim(X)
$.

\medskip\noindent
Now we explain joint work with Franccesco
Exposito from Padova.

First, a conjecture (by E-M).
Suppose $X$ is a compact complex manifold
with $\text{kod}(X)<\dim(X)$. If $\Omega_X^1$
is (weakly) semiample, then
$|\pi_1(X)|=\infty$.

The evidence is:
\begin{enumerate}
\item If $\dim X=1,2$, holds because of
classification.
\item If $X$ is projective, then the
conjecture follows from
a result of Hörng, which states that
$\Omega_X^1$ is semiample,
then there exists $X'\to X$ finite étale
cover
such that $X'\sim_{\text{bitat}}Y\times A$
where $A$
is an abelian variety and  $\dim
Y=\text{kod}(X)$.
\item If $X$ is Kähler, the conjecture (E-M)
follows
from a conjecture of Wu-Zheng from 2002,
which says that
if  $\Omega^1_X$ is (NEF) semiample (which
implies NEF), then
there exists $X' \to X$ finite étale that
admits a smooth fibration
$X' \to Y$ in complex tori with $\dim
Y=\text{kod}X$.
In this case the fundamental groups of the
tori are injected
in the fundamental group of $X$, i.e.,
$\pi_1(F) \hookrightarrow \pi_1(X)$.
(It appears that this is proved via the Serre
homotopy
long exact sequence, though the Kähler
hypothesis may needed, and
it's not clear how…)
This makes $|\pi_1(X)|=\infty$.

\item If $X$ is a compact complex manifold
and $\text{kod}(X)=0$,
and $\Omega_X^1$ is (weakly) semiample, then
$|\pi_1(X)|=\infty$.
\end{enumerate}

\noindent

\begin{definition}
\label{definition-}
Let $E$ be a vector bundle on $X$, a compact
complex manifold.
Recall the construction
$$
\xymatrix{
\mathbb{P}(E)\ar[d]^{\pi}\ar@{-}[r]&\pi^*
E\ar@{->>}[r]
&\mathcal{O}_{\mathbb{P}(E)}(1)\\
X\ar@{-}[r]&E
}
$$
$$
\pi_*(\mathcal{O}_{\mathbb{P}(E)}(m))=
\text{Sym}^m(E).
$$
\begin{enumerate}
\item $E$ is {\it strongly semiample} if
there exists $m>0$
such that $\text{Sym}^mE$ is globally
generated.
\item $E$ is {\it weakly semiample} if
$\mathcal{O}_{\mathbb{P}(E)}(1)$ is semiample
line bundle
(i.e. $\mathcal{O}_{\mathbb{P}(E)}(m)$ is
globally generated
for some $m\gg 0$).
\end{enumerate}
\end{definition}

\noindent
Notice that these two definitions are not
equivalent.
Strongly implies weakly, but not conversely.
To see that strongly implies weakly, consider
the pullback
$$
\xymatrix{
\pi^*\text{Sym}^nE=\text{Sym}^n\pi^*E\ar[r]
\ar[d]&\mathbb{P}(E)\ar[d]^{\pi}\\
\text{Sym}^mE\ar[r]&X
}
$$
then
 $$
\text{Sym}^m\pi^* E
\to\!\!\!\!\!\to\mathcal{O}
_{\mathbb{P}(E)}(m).
$$
To see that weakly does not imply strongly,
suppose $X$ is a smooth projective
surface with $\text{kod}(X)=0$.
Consider $E=\Omega_X^1$ semiample. We have
that $\Omega_X^1$
is strongly semiample if and only if $X$ is
an abelian surface.
Also, if $X$ is a hyperelliptic surface, then
$\Omega_X^1$
is weakly semiample.
This means that $X$ hyperelleptic surface is
the counterexample.

\begin{theorem}[E-M]
\label{theorem-em}
Suppose that $X$ is compact complex and that
$\text{kod}(X)=0$. Then
\begin{enumerate}
\item if $\Omega_X^1$ is strongly semiample
then
$X$ is a parallelalizable complex manifold.
\item If $\Omega_X^1$ is weakly semiample,
then $X$ is a finite étale Galois quotient
$X\cong P/G$
of a parallelizable manifold $P$.
\end{enumerate}
\end{theorem}

\begin{definition}
\label{definition-paralellizable}
$X$ is {\it parallelizable} is
$\Omega_X^1\cong\mathcal{O}_X$.
\end{definition}

\noindent
A $X$ is a compact complex manifold is
parallelizable if and only if
$X \cong H/\Gamma$ for $H$ a complex Lie
group and $\Gamma\subset H$
a discrete subgroup.
Further, it is Kähler if and only if $H$ is a
complex torus.

\section{Hyperkahler manifolds with 
finite group of birational automorphisms}
\label{section-hk-misha-26}

\noindent
Misha Verbitsky, 20 May, 2026. IMPA.

\medskip\noindent {\bf Abstract.} Seja $M$ uma
coletânea hyperkahler, $b_2(M)\ge 6$, e $D$ seu
espaço de deformação, $I\in D$ uma estrutura
complexa e $H_I$ a rede Hodge de $(M,I)$, ou seja,
o conjunto de todas as classes de cohomologia
inteiras do tipo $(1,1)$. Usando o teorema de
Vinberg sobre redes geradas por reflexões,
provamos que há apenas um número finito de
classes de isomorfismo de redes de Hodge $H_I$
de grau $\ge 3$, de modo que $(M,I)$ seja projetivo
e tenha um grupo finito de automorfismos
birracionais. Usamos essa observação para
mostrar que qualquer família não trivial de
deformações projetivas de $M$ tem um conjunto
denso de estruturas complexas com um grupo
infinito de automorfismos birracionais. Este
é um trabalho conjunto com Ekaterina Amerik e
Andrey Soldatenkov.  


\medskip\noindent
{\bf Upshot.}
The main result we shall prove is the
following.
Let $M$ be a K3 surface and let $\mathcal{H}$
be its polarized moduli space.
There exist finitely many points
$z_1,\dots,z_n \in \mathcal{H}$
such that for any
$m \in \mathcal{H}\setminus \bigcup_{i=1}^n \{z_i\}$
with $\rho(M)\leq 3$, one has
$\lvert \operatorname{Aut}(M)\rvert=\infty$.

In plain language: for most members of the
family, once the Picard rank is small enough,
the surface has infinitely many
automorphisms in ``almost every case'';
only finitely many special
parameter values have infinite automorphism
group.

My favorite idea from the talk is the
following, which I think is called
Noether-Lefschetz result:

\begin{quotation}
  ``Torelli tells us that there are roughly as many
K3 surfaces as Hodge decompositions. Fix a
polarization $p \in H^2(M,\mathbb{Z})$ with
$p^2>0$. There is a group acting on
$\mathbb{P}\operatorname{er}$, namely
$\operatorname{SO}(H^2(M,\mathbb{Z}),p)$, and the
quotient is quasiprojective. This is precisely the
moduli space of polarized K3 surfaces.

\noindent
Now suppose you have a sublattice
$\Lambda \subset H^2(M,\mathbb{Z})$. The
Noether-Lefschetz component is
\[
H_{\Lambda}
\;=\;
\{\,M \in \mathbb{P}\operatorname{er}/\operatorname{SO}
\mid \Lambda \subset H^{1,1}(M)\,\}.
\]
The classes in such quotient correspond
with the K3 that have finite automorphism
group!"
\end{quotation}


\medskip\noindent
The talk starts reviewing several notions
and results like:

\begin{theorem}
A K3 surface $M$ is projective if and only if
there exists a class $\lambda \in H^{1,1}(M,\mathbb{Z})$
such that $\int_M \lambda \wedge \lambda > 0$.
\end{theorem}

\begin{remark}
From the Riemann-Roch formula it follows that
any $(-2)$-class $\lambda$ is either effective or
its negative $-\lambda$ is effective, that is,
equal to the fundamental class of a curve.
\end{remark}

\begin{remark}
This implies that a class in the orthogonal
complement of a $(-2)$-class cannot be Kähler.
\end{remark}

\noindent
The follwing theorem uses the Demailly-Paun
characterization of Kählerness:

\begin{theorem}
Let $M$ be a K3 surface, and $S \subset
H^{1,1}(M,\mathbb{Z})$ the set of all
$(-2)$-classes. Consider the decomposition of
$\operatorname{Pos}(M) \setminus
\bigcup_{s\in S} s^\perp$ into connected
components. Then the Kähler cone of $M$ is one
of the connected components.
\end{theorem}

\begin{theorem}
Let $M$ be a K3 surface. The group
$\operatorname{Aut}(M)$ of holomorphic
isometries of $M$ is naturally identified with
the subgroup of $O(H^2(M,\mathbb{Z}))$ which
preserves the Kähler cone.
\end{theorem}

\begin{definition}
A group is called virtually free abelian if it
contains a free abelian subgroup of finite
index.
\end{definition}

\begin{proposition}[Oguiso]
Let $M$ be a K3 surface, and $H^{1,1}(M,\mathbb{Z})$
its Neron-Severi lattice. Then the following
alternatives occur.
\begin{enumerate}
\item[(a)] If the intersection form on
$H^{1,1}(M,\mathbb{Z})$ is negative definite, then
the kernel of $\Psi: \operatorname{Aut}(M)\to
\operatorname{Aut}(H^{2,0}(M))=\mathbb{C}^*$ is finite,
and the image is virtually free abelian.
\item[(b)] If the intersection form on
$H^{1,1}(M,\mathbb{Z})$ is degenerate, the image and
the kernel of the map $\Psi$ are virtually free
abelian.
\item[(c)] If $M$ is projective, the image of $\Psi$
is finite.
\end{enumerate}
\end{proposition}

\begin{lemma}
Let $M$ be a projective K3 surface with Picard
rank $1$. Then $\operatorname{Aut}(M)$ is finite.
\end{lemma}

\begin{proof}
Indeed, for a projective K3, the map
$\operatorname{Aut}(M)\to O(H^{1,1}(M,\mathbb{Z}))$
has finite kernel (Proposition 1).
\end{proof}

\begin{definition}
Let $(\Lambda,q)$ be a quadratic lattice, that is,
$\mathbb{Z}^n$ equipped with an integer-valued
quadratic form. We say that $(\Lambda,q)$
represents $n\in \mathbb{Z}$ if there exists
$x\in \Lambda$ such that $q(x,x)=n$.
\end{definition}

\begin{lemma}
Let $M$ be a projective K3 surface with Picard
rank $2$. Then $\operatorname{Aut}(M)$ is finite if
its Neron-Severi lattice represents zero, and
contains an infinite cyclic subgroup of finite
index otherwise.
\end{lemma}

\begin{proof}
Since $M$ is projective, $\operatorname{Aut}(M)$ is
finite if and only if its group $\operatorname{Aut}_{\Omega}(M)\subset
\operatorname{SO}(H^{1,1}(M,\mathbb{Z}))$ of symplectic
automorphisms is finite (Proposition 1). The group
of isometries of a quadratic lattice $(\Lambda,q)$
of signature $(1,1)$ is finite if $(\Lambda,q)$
represents zero. It contains an infinite cyclic
subgroup of finite index otherwise, by Pell's
classification of solutions of an equation
$x^2-ay^2=1$.
\end{proof}

\noindent
It looks like a ``divisorial
reflection'' is just a reflection
w.r.t. a  $(-2)$-curve.

\begin{definition}
Let $v \in H^{1,1}(M,\mathbb{Z})$ be a $(-2)$-class.
Consider the map
\[
x \mapsto x-\frac{2(x,v)}{(v,v)}v,
\]
where $(\cdot,\cdot)$ is the intersection form.
This map is called the divisorial reflection.
\end{definition}

\begin{lemma}
The group generated by divisorial reflections acts
transitively on the set of connected components of
\[
\operatorname{Pos}(M)\setminus \bigcup_{s\in S}s^\perp,
\]
where $S$ is the set of all $(-2)$-classes.
\end{lemma}

\noindent
We denote by $\operatorname{Ref}(M)$ the subgroup
of $\operatorname{SO}(H^2(M,\mathbb{Z}))$
generated by divisorial reflections.

\noindent
Clearly, it is a normal subgroup in the group of
all Hodge isometries of $H^2(M,\mathbb{Z})$.

\begin{lemma}
The group $\operatorname{Aut}_{\Omega}(M)$ of
symplectic automorphisms of a K3 surface is
determined, up to an isomorphism, by its
Neron-Severi quadratic lattice.
\end{lemma}

\begin{proof}
It is identified with the subgroup
$\operatorname{St}(\operatorname{Kah}(M))$ of
$\operatorname{SO}(H^{1,1}(M,\mathbb{Z}))$
preserving the Kähler cone. However,
$\operatorname{Ref}(M)$ acts on
$\operatorname{SO}(H^2(M,\mathbb{Z}))$ by Hodge
isometries, transitively on all connected
components of
$\operatorname{Pos}(M)\setminus
\bigcup_{s\in S}s^\perp$.
Therefore, the action of $\operatorname{Ref}(M)$ is
transitive on all groups of form $\operatorname{St}(A)$,
where $A$ is a connected component in
$\operatorname{Pos}(M)\setminus
\bigcup_{s\in S}s^\perp$, and $\operatorname{St}(A)$
its stabilizer.
\end{proof}

\begin{remark}
The group $O(m,n)$, $m,n > 0$, has $4$ connected
components. We denote the connected component of $1$
by $SO^+(m,n)$. We call a vector $v$ positive if its
square is positive.
\end{remark}

\begin{definition}
Let $V$ be a vector space with quadratic form $q$
of signature $(1,n)$, $\operatorname{Pos}(V)=\{x\in
V \mid q(x,x)>0\}$ its positive cone, and
$\mathbb{P}V$ projectivization of $\operatorname{Pos}(V)$.
Denote by $g$ any $SO(V)$-invariant Riemannian
structure on $\mathbb{P}V$. Then $(\mathbb{P}V,g)$ is
called hyperbolic space, and the group $SO^+(V)$ the
group of oriented hyperbolic isometries.
\end{definition}

\begin{example}
For any K3 surface, $\mathbb{P}\operatorname{Pos}(M)$
is a hyperbolic space, and $\operatorname{Aut}(M)$
acts on this space by isometries.
\end{example}

\begin{theorem}[Sterk, 1985]
Let $M$ be a projective K3 surface. Then
$\operatorname{Aut}(M)$ acts with finitely many orbits
on the set of the walls of the Kähler cone.
\end{theorem}

\noindent
The proof may be found in a book by Kneser.

\medskip\noindent
When $M$ is a Calabi-Yau manifold, Sterk's
theorem is known as the Morrison-Kawamata
conjecture, formulated by Morrison in 1993,
and proven by Kawamata for elliptic Calabi-Yau
manifolds in 1997. It was proven by B. Totaro
for log-Calabi-Yau surfaces (2010), and by
Amerik-V. for hyperkähler manifolds (2015).

\begin{remark}
This is an AI-generated explanatory remark.
The ample cone is usually described by ample line
bundles, so its natural generators have integral
first Chern classes. Rational classes appear here
because one often passes to the rational span of
the cone when intersecting with the Kähler cone:
on a projective K3 surface, ample classes are
exactly the rational Kähler classes.
\end{remark}

\begin{lemma}
Let $\operatorname{Amp}(M):=\operatorname{Kah}(M)\cap H^{1,1}(M,\mathbb{Q})$
denote the ample cone of a projective K3
surface, and $\operatorname{Amp}_{\mathbb{R}}(M):=
\operatorname{Kah}(M)\cap H^{1,1}(M,\mathbb{Q})\otimes_{\mathbb{Q}}\mathbb{R}$
denote the corresponding cone in
$H^{1,1}(M,\mathbb{Q})\otimes_{\mathbb{Q}}\mathbb{R}$.
Consider its projectivization
$\mathbb{P}\operatorname{Amp}_{\mathbb{R}}(M)$ as a
subset in the hyperbolic space $\mathbb{P}\operatorname{Pos}(M)$.
Then $\mathbb{P}\operatorname{Amp}_{\mathbb{R}}(M)$
has finite volume if and only if
$\operatorname{Aut}(M)$ is finite.
\end{lemma}

\begin{proof}
By construction,
$\mathbb{P}\operatorname{Amp}_{\mathbb{R}}(M)$ is a
hyperbolic polyhedron (with some vertices on
absolute). If $\operatorname{Aut}(M)$ is finite, this
polyhedron has finite number of faces, hence its
volume is finite. Conversely, if $\operatorname{Aut}(M)$
is infinite, $\mathbb{P}\operatorname{Amp}_{\mathbb{R}}(M)$
contains infinite number of copies of its fundamental
domain, hence its volume is infinite.
\end{proof}

\begin{theorem}[Vinberg]
Let $M$ be a projective K3 surface of
Picard rank $\geq 3$. Then $\operatorname{Aut}(M)$
is finite if and only if the reflection group
$\operatorname{Ref}(M)$ has finite index in
$\operatorname{SO}(H^{1,1}(M,\mathbb{Z}))$.
\end{theorem}

\begin{proof}
Step 1: If $\operatorname{Aut}(M)$ is finite, the
ample cone has finite volume, hence the action of
$\operatorname{Ref}(M)$ on the hyperbolic space
$\mathbb{P}\operatorname{Amp}_{\mathbb{R}}(M)$ has
finite volume. This implies that $\operatorname{Ref}(M)$
is a lattice subgroup in
$\operatorname{SO}(H^{1,1}(M,\mathbb{R}))$, hence has
finite index in $\operatorname{SO}(H^{1,1}(M,\mathbb{Z}))$.

Step 2: Conversely, assume that $\operatorname{Ref}(M)$
has finite index in $\operatorname{SO}(H^{1,1}(M,\mathbb{Z}))$.
Then its Coxeter polyhedron is finite (Vinberg).
However, $\mathbb{P}\operatorname{Amp}_{\mathbb{R}}(M)$
is its Coxeter polyhedron.
\end{proof}

\begin{definition}
\label{definition-reflective}
We call a quadratic lattice with
$\operatorname{Ref}(\Lambda)$ of finite index in
$\operatorname{SO}(\Lambda)$ reflective.
\end{definition}

\noindent
Torelli tells us that there are roughly as many
K3 surfaces as Hodge decompositions. Fix a
polarization $p \in H^2(M,\mathbb{Z})$ with
$p^2>0$. There is a group acting on
$\mathbb{P}\operatorname{Per}$, namely
$\operatorname{SO}(H^2(M,\mathbb{Z}),p)$, and the
quotient is quasiprojective. This is precisely the
moduli space of polarized K3 surfaces.

\noindent
Now suppose you have a sublattice
$\Lambda \subset H^2(M,\mathbb{Z})$. The
Noether-Lefschetz component is
\[
H_{\Lambda}
\;=\;
\{\,M \in \mathbb{P}\operatorname{Per}/\operatorname{SO}
\mid \Lambda \subset H^{1,1}(M)\,\}.
\]
The classes in such quotient correspond
with the K3 that have finite automorphism
group!

\begin{lemma}
\label{lemma-noether-lefschetz-components}
Let $H$ be the moduli space of polarized K3
surfaces. Then there is a finite set of
quasiprojective subvarieties
$Z_1,\dots,Z_n \subset H$ such that for any K3
surface $M$ outside of $\bigcup_i Z_i$ with Picard
rank $\geq 3$, the group $\operatorname{Aut}(M)$ is
infinite.
\end{lemma}

\begin{proof}
Step 1: Fix a polarization on $M$, that is, an
integer vector $p \in H^2(M,\mathbb{Z})$ such that
its square is positive. By Torelli theorem,
$H=\mathbb{P}\operatorname{Per}/\operatorname{SO}(H^2(M,\mathbb{Z}),p)$,
where $\operatorname{Per}$ is the space of all Hodge
structures on $H^2(M)$ such that $p$ has type
(1,1), and $\operatorname{SO}(H^2(M,\mathbb{Z}),p)$
the group of isometries preserving $p$. For any
sublattice $\Lambda \subset H^2(M,\mathbb{Z})$
containing $p$, the corresponding
Noether-Lefschetz component $H_{\Lambda}\subset H$
is the image of all $I \in \operatorname{Per}$ such
that $H^{1,1}(M,I)\subset \Lambda$. By Baily-Borel
theorem, $H_{\Lambda}$ is a quasiprojective
subvariety in $H$.

Step 2: Let $Z_i$ be the set of Noether-Lefschetz
components associated with reflective lattices
$\Lambda \subset H^2(M,\mathbb{Z})$. Clearly, any
$M$ which has Picard rank $\geq 3$ and is not in
$Z_i$ has infinite automorphism group.
\end{proof}

\begin{theorem}[K. Oguiso, 2003]
\label{theorem-oguiso-2003}
Let $M_t$ be a non-trivial quasiprojective family
of K3 surfaces over a curve $C$, with general
member of Picard rank $\geq 3$. Then
$|\operatorname{Aut}(M_t)|=\infty$ for a dense set
of $t \in C$.
\end{theorem}

\begin{proof}
Remove from $C$ all proper subsets which belong to
the Noether-Lefschetz components $Z_i$. If
$C \not\subset Z_j$, we are done. Otherwise, there
is a dense subset of points in $C$ where the Picard
rank jumps, and these points are not in $Z_i$.
\end{proof}

\begin{remark}
\label{remark-oguiso-lower-picard}
This theorem is also valid for families with
general member of Picard rank $\geq 1$, but the
proof for lower Picard values is more elementary.
\end{remark}

\begin{remark}
\label{remark-oguiso-vinberg}
Oguiso's proof does not use Vinberg's theorem, he
uses explicit lattice-theoretic arguments.
\end{remark}

\begin{theorem}[Denisi, Onorati, Rizzo, Viktorova]
\label{theorem-denisi-onorati-rizzo-viktorova}
Let $M_t$ be a non-trivial quasiprojective family
of hyperkähler manifolds of known type over a
curve $C$. Then $|\operatorname{Bir}(M_t)|=\infty$
for a dense set of $t \in C$, where
$\operatorname{Bir}(M_t)$ denotes the group of
birational automorphisms.
\end{theorem}

\begin{remark}
\label{remark-dorv-lattices}
Their proof relies on known properties of the
cohomology lattices of hyperkähler manifolds.
\end{remark}

\begin{theorem}[Amerik, Soldatenkov, V.]
\label{theorem-amerik-soldatenkov-v}
This theorem holds for any family of hyperkähler
manifolds.
\end{theorem}

\begin{remark}
\label{remark-asv-vinberg}
Our proof is based on Vinberg's theorem, in fact,
it is pretty much identical to the proof for K3
surfaces that I gave.
\end{remark}

\section{Multiplicity of the fibers in
Lagrangian fibrations}
\label{section-multiplicity-fibers-
lagrangian-fibrations}

\noindent
Misha Verbitsky.
Ter 02 jun 2026, 15:30.

\medskip\noindent
Seja \(M\) um coletor hyperkahler de
holonomia m\'axima e
\(\pi : M \to \mathbf{P}^n\) uma fibração
lagrangiana. Uma fibra \(F\) de \(\pi\) \'
e m\'ultipla de multiplicidade \(p\) se um
pequeno disco transversal a \(F\) encontrar
uma fibra geral de \(\pi\) em \(p > 1\)
pontos. Provarei que as fibras m\'ultiplas
n\~ao ocorrem em codimens\~ao 1 se e
somente se o pullback de uma se\c{c}\~ao de
hiperplano for primitivo (indivis\'ivel) na
cohomologia de \(M\). Em seguida,
explicarei como usar essa observa\c{c}\~ao
para provar que n\~ao h\'a fibras m\'ultiplas
em codimens\~ao 1. Os argumentos se
baseiam em uma vers\~ao forte do
desaparecimento de Kodaira-Nakano, v\'alido
para feixes de linhas com curvatura
semipositiva, e na estabilidade do feixe
tangente. Este \'{e} um trabalho conjunto
com Ljudmila Kamenova.

\medskip\noindent
It's a mystery why any elliptic fibration
over a K3 surface has no multiple fibers.
We will explain how to generalize this result
to higher dimensions. By some
theorem, elliptic fibrations shall
be Lagrangian fibrations.

\medskip\noindent

\begin{theorem}[Calabi--Yau]
\label{theorem-calabi-yau}
A compact, K\"ahler, holomorphically
symplectic manifold admits a unique
hyperk\"ahler metric in any K\"ahler class.
\end{theorem}

\begin{definition}
\label{definition-hyperkahler-manifold}
For the rest of this talk, a hyperk\"ahler
manifold is a compact, K\"ahler,
holomorphically symplectic manifold.
\end{definition}

\begin{definition}
\label{definition-maximal-holonomy}
A hyperk\"ahler manifold \(M\) is called
maximal holonomy, or IHS, if
\(\pi_1(M)=0\) and \(H^{2,0}(M)=\mathbf{C}\).
\end{definition}

\medskip\noindent
\textbf{Bogomolov's decomposition.}
Any hyperk\"ahler manifold admits a finite
covering which is a product of a torus and
several hyperk\"ahler manifolds of maximal
holonomy.

\medskip\noindent
Further on, all hyperk\"ahler manifolds are
assumed to be of maximal holonomy.

\begin{theorem}[Matsushita]
\label{theorem-matsushita-lagrangian}
Let \(M\) be a hyperk\"ahler manifold of
maximal holonomy, and let
\(\pi : M \to X\) be a surjective
holomorphic map, with \(0 < \dim X < \dim M\).
Then \(\pi\) is a Lagrangian fibration
(that is, it has holomorphic Lagrangian
fibers).
\end{theorem}

\begin{theorem}[Hwang]
\label{theorem-hwang-smooth-base}
In these assumptions, \(X\) is biholomorphic
to \(\mathbf{C}P^n\) when it is smooth.
\end{theorem}

\medskip\noindent
The conjecture is that \(X\) is biholomorphic
to \(\mathbf{C}P^n\) when it is normal.

\begin{theorem}[Matsushita]
\label{theorem-matsushita-normal-base}
Let \(M\) be a hyperk\"ahler manifold of
maximal holonomy, and let
\(\pi : M \to X\) be a Lagrangian fibration,
with \(X\) normal. Then
\(H^*(X,\mathbf{Q}) \cong H^*(\mathbf{C}P^n,
\mathbf{Q})\).
\end{theorem}

\begin{remark}
\label{remark-general-fibers-abelian}
General fibers of \(\pi\) are Abelian
varieties (projective complex tori), by
Arnold-Liouville. Conversely, as shown by
Hwang-Weiss, any Lagrangian complex torus in
\(M\) is a fiber of a Lagrangian fibration.
\end{remark}

\begin{remark}
\label{remark-fibration}
A fibration \(\pi : M \to X\) is a proper
surjective holomorphic map of complex
manifolds. We shall always tacitly assume
that the fibration is proper and
equidimensional, that is, the irreducible
components of all fibers have dimension
\(\dim M - \dim X\). A special fiber of
\(\pi\) is a fiber containing a critical
point.
\end{remark}

\begin{definition}
\label{definition-multiplicity-special-fiber}
Let \(Z\) be an irreducible component of a
special fiber of a fibration
\(\pi : M \to X\), and \(z=\pi(Z)\). The
multiplicity of \(Z\) is the rank of
\(\pi^*(\mathcal{O}_X/\mathfrak{m}_z)\) in a
general point of \(Z\), where
\(\mathfrak{m}_z \subset \mathcal{O}_X\) is
the maximal ideal of \(z\).
\end{definition}

\begin{remark}
\label{remark-multiplicity-one-dimensional}
Suppose that \(\dim X=1\). Locally in \(M\),
around a general point of \(m \in Z\), there
is a coordinate system \(x_1,\ldots,x_n\),
such that \(\pi(x_1,\ldots,x_n)=x_1^d\).
In this case, \(d\) is multiplicity of
\(Z\).
\end{remark}

\begin{remark}
\label{remark-equivalent-definition}
Another definition of multiplicity is more
geometric, but it is equivalent to the one
given above; the equivalence is left as an
exercise.
\end{remark}

\begin{definition}
\label{definition-multiplicity-transversal-disk}
Let \(Z\) be an irreducible component of a
special fiber of a fibration
\(\pi : M \to X\), and \(z=\pi(Z)\).
Consider a small disc \(D \subset M\) of
dimension \(\dim X\) transversal to \(Z\)
and intersecting it in \(m\). The
multiplicity of \(Z\) in \(m \in M\) is the
number of intersection points between \(D\)
and a general fiber of \(\pi\) which is
sufficiently close to \(Z\).
\end{definition}

\begin{theorem}[Campana-Kamenova-Verbitsky]
\label{theorem-campana-kamenova-verbitsky}
Let \(\pi : M \to X\) be an Abelian
fibration (that is, a fibration with general
fiber a complex torus), and assume that
\(M\) is K\"ahler. Let \(\mu_i\) be the
multiplicity of irreducible components of
\(\pi^{-1}(z)\) for some \(z \in X\), and
\(\gcd(\mu_i)\) their greatest common
divisor. Then \(\gcd(\mu_i)=\min \mu_i\).
\end{theorem}

\begin{definition}
\label{definition-discriminant}
Let \(\pi : M \to X\) be a surjective
holomorphic map of complex manifolds, and
\(D \subset X\) its set of critical values,
which is known as the discriminant, or the
discriminant divisor (it has codimension
1). We say that \(\pi\) has no multiple
fibers in codimension 1 if for a general
point \(x \in D\), the fiber \(\pi^{-1}(x)\)
has a component with multiplicity 1.
\end{definition}

\begin{theorem}
\label{theorem-elliptic-k3-no-multiple-fibers}
Let \(\pi : M \to X\) be an elliptic
fibration on a K3 surface. Then \(\pi\) has
no multiple fibers.
\end{theorem}

\medskip\noindent
The main result today is a generalization of
this theorem.

\begin{theorem}
\label{theorem-main-no-multiple-fibers}
Let \(\pi : M \to X\) be a Lagrangian
fibration on a hyperk\"ahler manifold. Then
\(\pi\) has no multiple fibers in
codimension 1.
\end{theorem}

\begin{lemma}
\label{lemma-h2-torsion-free}
Let \(M\) be a hyperk\"ahler manifold of
maximal holonomy. Then \(H^2(M)\) is
torsion-free.
\end{lemma}

\begin{proof}
The universal coefficients formula gives the
exact sequence:
$$
0 \to \operatorname{Ext}_{\mathbb{Z}}^1
\bigl(H_1(M;\mathbb{Z}),\mathbb{Z}\bigr)
\to H^2(M;\mathbb{Z}) \to
\operatorname{Hom}_{\mathbb{Z}}
\bigl(H_2(M;\mathbb{Z}),\mathbb{Z}\bigr)
\to 0.
$$
Since \(H_1(M,\mathbb{Z})=0\) for a maximal
holonomy hyperk\"ahler manifold, this gives
an isomorphism
\[
H^2(M,\mathbb{Z}) =
\operatorname{Hom}_{\mathbb{Z}}
\bigl(H_2(M;\mathbb{Z}),\mathbb{Z}\bigr),
\]
hence the torsion vanishes.
\end{proof}

\begin{lemma}
\label{lemma-primitive-classes-lagrangian}
Let \(\pi : M \to \mathbf{C}P^n\) be a
Lagrangian fibration on a hyperk\"ahler
manifold of maximal holonomy, and
let \(H \subset \mathbf{C}P^n\) be a
hyperplane section. Then the following
assertions are equivalent.
\begin{itemize}
\item The homology class of
\(\pi^{-1}(H)\) is primitive.
\item The map \(\pi\) has no multiple
fibers in codimension 1.
\end{itemize}
\end{lemma}

\begin{proof}
Step 1: The implication (ii) \(\Rightarrow\)
(i) is proven later today using the ETMDPS
vanishing theorem. The converse implication:
let \(D_1\) be an irreducible component of
the discriminant \(D\); assume it has
multiplicity \(\mu\). We need to show that
the preimage of the hyperplane section \(H\)
is not primitive.

Step 2: Arguing ad absurdum, assume that
\(D_1\) is homologous to \(k[H]\), but
\([\pi^{-1}(H)]\) is primitive in
\(H^2(M,\mathbb{Z})\). Since
\([D_1] = k[H]\), \(\pi^{-1}(D_1)\) is the
zero set of a section
\(s \in H^0(\mathbf{C}P^n,\mathcal{O}(k))\).
The pullback of this section has multiplicity
\(\mu\), since \(\pi^*(\mathcal{O}(1))\) is
primitive, \(k\) is divisible by \(\mu\).
The universal coefficients formula implies
that \(H^2(M,\mathbb{Z})\) is torsion-free.
Therefore, \(\pi^*s = s_0^\mu\), where
\(s_0 \in H^0(M,\pi^*\mathcal{O}(k/\mu))\).
Since \(\pi^*\mathcal{O}(k/\mu)\) is
trivial on fibers of \(\pi\), the section
\(s_0\) is a pullback of a section
\(s_1 \in H^0(\mathbf{C}P^n,\mathcal{O}(k/\mu))\).
This is impossible, because \(s\) and hence
\(s_1\) vanishes on \(D_1\), which has
degree \(k\).
\end{proof}

\begin{theorem}[Fujiki]
\label{theorem-fujiki}
Let \(\eta \in H^2(M)\), and \(\dim M = 2n\),
where \(M\) is hyperk\"ahler. Then
\[
\int_M \eta^{2n} = cq(\eta,\eta)^n,
\]
for some primitive integer quadratic form
\(q\) on \(H^2(M,\mathbb{Z})\), and
\(c>0\) a rational number.
\end{theorem}

\begin{definition}
\label{definition-bbf-form}
This form is called the Bogomolov-Beauville-
Fujiki form. It is defined by the Fujiki's
relation uniquely, up to a sign. The sign is
determined from the following formula
\[
\lambda q(\eta,\eta)=
\int_X \eta \wedge \eta \wedge \Omega^{n-1}
\wedge \overline{\Omega}^{n-1}
- \frac{n-1}{2n}
\left(
\int_X \eta \wedge \Omega^{n-1}
\wedge \overline{\Omega}^{n}
\right)
\left(
\int_X \eta \wedge \Omega^{n}
\wedge \overline{\Omega}^{n-1}
\right)
\]
where \(\Omega\) is the holomorphic symplectic
form, and \(\lambda>0\).
\end{definition}

\begin{remark}
\label{remark-bbf-signature}
\(q\) has signature \((3,b_2-3)\). It is
negative definite on primitive forms, and
positive definite on \(\langle \Omega,
\overline{\Omega}, \omega \rangle\), where
\(\omega\) is a K\"ahler form.
\end{remark}

\begin{definition}
\label{definition-holomorphic-euler}
Let \(B\) be a holomorphic vector bundle (or
a coherent sheaf). The holomorphic Euler
characteristic is
\[
\chi(B) := \sum_i (-1)^i H^i(M,B).
\]
\end{definition}

\begin{theorem}[Riemann-Roch-Hirzebruch]
\label{theorem-rriemann-roch-hirzebruch}
Let \(M\) be a compact complex manifold, and
\(B\) a holomorphic vector bundle. Then
\(\chi(B)\) can be expressed through the
the Chern classes of \(TM\) and \(B\),
\[
\chi(B)=\int_M td(TM)\wedge ch(B)
\]
where \(td\) is the the Todd polynomial on
Chern classes of \(TM\),
\[
td(M)=1+\frac{1}{2}c_1+\frac{1}{12}(c_1^2+c_2)
+\frac{1}{24}c_1c_2+\frac{1}{720}
(-c_1^4+4c_1^2c_2+c_1c_3+3c_2^2-c_4)+\ldots
\]
and \(ch(B)\) its Chern character,
\[
ch(B)=1+c_1+\frac{1}{2}(c_1^2-2c_2)
+\frac{1}{6}(c_1^3-3c_1c_2+3c_3)+\ldots
\]
\end{theorem}

\noindent
The idea of the following theorem
is to look at $\text{SO}(H^2)$.
The function $\chi(L)$ is an isometry
of BBF form, which makes it a polynomial
on $q(v,v)$. Also, isometries of $H^2$
are correspond with automorphisms of $M$.

\begin{theorem}[Huybrechts]
\label{theorem-huybrechts}
Let \(M\) be a hyperk\"ahler manifold,
\(\dim_{\mathbf{C}} M = 2n\), and \(L\) a
holomorphic line bundle. Then
\[
\chi(L)=\sum a_i q(c_1(L))^i,
\]
where the coefficients \(a_i\) are constants
depending on the topology of \(M\).
\end{theorem}

\begin{proof}
Step 1: Let \(A^*\) be the subalgebra in
cohomology generated by \(H^2(M)\). Then
\[
A^{2i} \cong \operatorname{Sym}^i(H^2(M))
\]
up to the middle degree, and
\[
A^{n+i} \cong \operatorname{Sym}^{n-i}
\bigl(H^2(M)\bigr);
\]
there is an \(O(H^2(M))\)-action on
cohomology, and the multiplication is
\(O(H^2(M))\)-invariant (V., 1995).

Step 2: All Chern classes of \(TM\) are
\(O(H^2(M))\)-invariant, but there is only
one (up to a constant multiplier)
\(O(H^2(M))\)-invariant functional on
\(\operatorname{Sym}^{2i}(H^2(M))\). On the
class \(\eta^{2i} \in H^{4i}(M)\) this
functional takes value \(q(\eta,\eta)^i\).
Therefore, all \(L\)-dependent coefficients
in the Hirzebruch-Riemann-Roch formula for
\(\chi(L)\) are expressed through
\(q(c_1(L))\).
\end{proof}

\begin{lemma}[Corollary]
\label{corollary-huybrechts-chi}
Let \(L\) be a line bundle on a hyperk\"ahler
manifold \(M\), \(\dim_{\mathbf{C}} M = 2n\).
Assume that \(q(c_1(L))=0\). Then
\(\chi(L)=n+1\).
\end{lemma}

\begin{proof}
Indeed, \(\chi(L)=\chi(\mathcal{O}_M)=n+1\),
with the second equality implied by
Bochner's vanishing theorem.
\end{proof}

\begin{definition}
\label{definition-semipositive}
A real \((1,1)\)-form \(\eta\) on a complex
manifold \(M\) is called semipositive if
\(\eta(x,Ix)\geq 0\) for all real tangent
vectors \(x\).
\end{definition}

\medskip\noindent
The following theorem was rediscovered
several times during 1990-ies (Enoki,
Takegoshi, Morougane). Its most general form
(which we do not use) is due to Demailly,
Peternell and Schneider.

It is similar to Kodaira-Nakano
vanishing theorem: instead of having that
$H^i(L \otimes K)$ will vanish for ample $L$,
we have an estimate for how big it can be:

\begin{theorem}[ETMDPS]
\label{theorem-etmdps-vanishing}
Let \((M,I,\omega)\) be a compact K\"ahler
manifold, \(\dim_{\mathbf{C}} M = n\), \(K\)
its canonical bundle, and \(L\) a
holomorphic line bundle on \(M\) equipped
with a Hermitian metric \(h\). Assume that
the curvature \(\Theta\) of \(L\) is a
positive form on \(M\). Then the wedge
multiplication operator \(\eta \mapsto
\omega^i \wedge \eta\) induces a surjective
map
\[
H^0(\Omega^{n-i}_M \otimes L)
\xrightarrow{\ \omega^i\ }
H^i(K \otimes L).
\]
Here \(\omega\) is considered as an element
in \(H^1(\Omega^1_M)\), and multiplication by
\(\omega\) maps
\[
H^k(\Omega^{n-l}_M \otimes L)
to H^{k+1}(\Omega^{n-l+1}_M \otimes L).
\]
\end{theorem}

\begin{lemma}
\label{lemma-global-sections-vanish}
Let \(B\) be a vector bundle equipped with a
filtration \(0 = B_0 \subset B_1 \subset
\ldots \subset B_k = B\). Assume that
\(H^0(B_i/B_{i-1}) = 0\) for all \(i\).
Then \(H^0(B)=0\).
\end{lemma}

\begin{theorem}
\label{theorem-line-bundle-trivial-on-fibers}
Let \(M\) be a hyperk\"ahler manifold
admitting a Lagrangian fibration
\(\pi : M \to X\), and \(H\) a line bundle on
\(X\). Let \(L\) be a line bundle such that
\(L^{\otimes k} = \pi^*H\). Then \(L\) is
trivial on all smooth fibers of \(\pi\).
\end{theorem}

\begin{proof}
Step 1: Let \(F\) be a smooth fiber of
\(\pi\), which is an abelian variety by
Arnol'd-Liouville. Then \(T^*M|_F\) is an
extension of a trivial bundle \(TF\) with
another trivial bundle \(NF = T^*F\). For
any non-trivial line bundle
\(L \in \operatorname{Pic}_0(F)\), we have
\(H^0(L\otimes TF)=0\) and
\(H^0(L\otimes NF)=0\), which implies that
\(H^0(L\otimes T^*M|_F)=0\). Similarly one
\(H^0(L\otimes \wedge^k T^*M|_F)=0\)
(Lemma 2).
Step 2: Unless \(L\) is trivial on \(F\), we
have \(H^0(L\otimes \wedge^k T^*M|_F)=0\),
which implies that \(H^0(L\otimes \wedge^*M)
=0\). By Enoki-Mourugane-Takegoshi-Demailly-
Peternell-Schneider theorem, this implies
that \(H^i(L)=0\), hence \(\chi(L)=0\),
contradicting the formula \(\chi(L)=n+1\).
\end{proof}

Consider a fibration \(M \to \mathbf{C}P^n\)
with fiber \(T^n\). This gives an exact
sequence
\[
0 \to TF \to TM|_F \to \pi^*T\mathbf{C}P^n|_F
\to 0.
\]
The point is that when we have a
decomposition into trivial bundles, it is
very easy to control the sections.
The usage appears to be that a torsion
bundle tensored with a trivial bundle cannot
have sections.

\begin{proposition}
\label{proposition-fiberwise-monodromy-line-bundle}
Let \(M\) be a hyperk\"ahler manifold
admitting a Lagrangian fibration
\(\pi : M \to X\), and \(H\) a line bundle on
\(X\). Let \(L\) be a line bundle such
that \(L^{\otimes k} = \pi^*H\). Then \(L\)
admits a connection \(\nabla\) which is flat
on each restriction \(L|_F\) to the fiber of
\(\pi\).
\end{proposition}

\begin{proof}
Choose a constant metric \(h^k\) on
\(L^{\otimes k}|_F = \mathcal{O}_F\) and let
\(h\) be its \(k\)-th root, which is a metric
on \(L|_F\). Since \(h^k\) is constant on
each fiber, its curvature vanishes, and the
Chern connection \(\nabla\) associated with
\(h\) is also flat.
\end{proof}

\begin{definition}
\label{definition-fiberwise-monodromy}
Fiberwise monodromy of \(L\) is its
monodromy on the fibers of \(\pi\).
\end{definition}

\begin{remark}
\label{remark-pullback-monodromy}
Clearly, \(L = \pi^*L_0\) if and only if the
fiberwise monodromy of \(L\) on each fiber
is trivial.
\end{remark}

\begin{remark}
\label{remark-flat-connection}
Since \(\mathcal{O}(1)\) is flat on
\(\mathbf{C}^n = \mathbf{C}P^n \setminus H\),
we also obtain a flat connection on
\(L|_{\pi^{-1}(\mathbf{C}P^n\setminus H)}\).
The main theorem would follow if we prove
that the monodromy of this flat connection is
trivial on \(\pi^{-1}(z)\), where
\(z \in D\) is a general point of the
discriminant.
\end{remark}

\section{G4 fluxes and Hodge theory}
\label{section-g4-fluxes-and-hodge-theory}

\noindent
Daniel Lopez, UFF.
GADEPS Focused Meeting June 30, 2026, IMPA.

\medskip\noindent
A flux $G$ is a primitive Hodge class,
i.e. an element of $H^{2,2}(X)
\cap H^4(X,\mathbb{Z})$.

\medskip\noindent
Conjecture: if $G$ is general,
 $|G|\gg \chi(X)/12$.

\medskip\noindent
Extra note:
$h^{3,1}$ is the dimension of the moduli
of fourfolds.

\noindent
The second method (the first
one is Griffiths basis)
uses Poincaré duality to pass
from (rational) Hodge classes
to (rational) homology classes and in
this language formulate hodge conjecture.
Namely, that any rational Hodge homology
class is algebraic.

\begin{example}
\label{example-fermat-sextic}
The {\it Fermat sextic fourfold} 
$X \subset \mathbb{P}^5$ is the hyupersurface
given by
$$
F(x) = X_0^6+….+x_5^6.
$$
It has a lot of symmetries;
its automorphism group
is given as a semidirect product
and has an action
of symmetric group on automorphisms.

The Griffiths basis is
$$
\omega_\beta=
\operatorname{Res}
\left(\frac{x^\beta \Omega}{F^{k+1}}\right)
\in H^{4-k,k}
$$
where $\beta$ is the multiindex
$x^\beta= x_0^{\beta_0}…x_5^{\beta_5}$.
This information allows us to compute
the Hodge numbers of $X$.

Shioda proved in 1979 that $X$ satisfies
the Hodge conjecture over $\mathbb{Q}$.
Braun, Fortin, Lopez, Villaflor, in 2024,
provided a basis for $H^{2,2}(X) \cap
H^4(X,\mathbb{Q})$ given by algebraic cycles
divided into 3 different types.
\end{example}

\begin{example}
\label{example-weighted-36-fermat}
Consider the {\it weighted 36 Fermat} 
given by
$$
F=x_0^{36}+x_1^{36}+x_2^{36}+x_3^4
+x_4^4+x_5^3=0
$$
in $\mathbb{P}_{1,1,1,9,12,12}$.

It's not known whether the Hodge conjecture
holds for this variety.
\end{example}



\end{document}

